\exercisesection{REGULAR LOCAL RINGS}

\begin{enumerate}
\def\labelenumi{\arabic{enumi})}
\item
  Let A be a ring, S a multiplicative subset of A. Show that one has \(\mathrm{dh}(\mathrm{S}^{-1}\mathrm{A}) \leqslant \mathrm{dh}(\mathrm{A})\) (A, X, p.138, def. 2). (Take resolutions of the \(S^{-1}A\) -modules by projective A-modules and localize.)
\item
  Let A be a Noetherian ring and B a faithfully flat A-algebra. Then we have \(\mathrm{dh(A)}\leqslant \mathrm{dh(B)}\).
\end{enumerate}

3) Let A be a Noetherian local ring. We propose to show that A is regular if and only if dh(A) is finite.

\begin{enumerate}
\def\labelenumi{\alph{enumi})}
\tightlist
\item
  Suppose dh(A) is finite. To show that A is regular, use A, X, p.208, exerc. 13.
\item
  Suppose A is regular. Let x be a completely secant family generating m\_A. Then, the Koszul complex associated with x is a finite resolution of A/m\_A by free A-modules. Then apply A, X, p.203, exerc. 16.
\end{enumerate}

\begin{enumerate}
\def\labelenumi{\arabic{enumi})}
\setcounter{enumi}{3}
\item
  Let A be a regular Noetherian local ring. Then we have \(\dim(A) = \mathrm{dh}(A)\) (To show that we have \(\mathrm{dh}(A) \leqslant \dim(A)\), use the Koszul complex associated with a system of coordinates. Then determine the groups \(\operatorname{Ext}_{A}^{i}(\kappa_{A}, \kappa_{A})\).)
\item
  Let A be a regular Noetherian local ring, and let p be a prime ideal of A. Then \(A_{p}\) is regular (use Exerc. 1 and Exerc. 3).
\item
  A Noetherian ring A is said to be regular if for every prime ideal p of A, the local ring \(A_{p}\) is regular.
\end{enumerate}

\begin{enumerate}
\def\labelenumi{\alph{enumi})}
\item
  Let A be a Noetherian ring. Then A is regular if and only if \(A_{m}\) is regular for every maximal ideal m of A (use Exercise 5).
\item
  Let A be a Noetherian ring. Show that A is regular of finite dimension if and only if dh(A) is finite (use Exercise 4).
\item
  Show that the ring described in Exercise 13, p.83 is regular of infinite dimension.
\item
  Let A be a regular Noetherian ring, and S a multiplicative subset. Show that \(\mathbf{S}^{-1}\mathbf{A}\) is regular.
\item
  Let A be a regular Noetherian ring. Show that A{[}X{]} is regular. (Reduce by localization to the case where dh(A) is finite. Use A, X, p. 143, th. 1, and b).
\end{enumerate}

\(^{1}\) For further details on this exercise, one may consult R. Stanley, The Hilbert function of a graded algebra, Advances in Math., 28 (1978), pp. 57–83.

\begin{enumerate}
\def\labelenumi{\arabic{enumi})}
\setcounter{enumi}{6}
\tightlist
\item
  Let I be a directed set admitting a countable cofinal subset, and let A be a ring (not necessarily commutative). We propose to show that every inductive limit, indexed by I, of projective (left) A-modules has projective dimension \(\leqslant 1\) (A, X, p.~134).
\end{enumerate}

\begin{enumerate}
\def\labelenumi{\alph{enumi})}
\tightlist
\item
  Let \(((\mathbf{Q}_i)_{i\in \mathbf{I}},(\varphi_{i,j})_{i < j})\) be an inductive system of projective A-modules and \(Q = \varinjlim_{I} Q_i\). To show that \(Q\) is of projective dimension \(\leqslant 1\), reduce to the case where \(I = \mathbf{N}\). We then have an exact sequence
\end{enumerate}

\[
0 \longrightarrow \bigoplus_ {n \in \mathbf {N}} Q _ {n} \stackrel {u} {\longrightarrow} \bigoplus_ {n \in \mathbf {N}} Q _ {n} \longrightarrow Q \longrightarrow 0
\]

where for every \(n \in \mathbf{N}\) and every \(x \in \mathbf{Q}_n\) , one has

\[
u (x) = \varphi_ {n, n + 1} (x) - x.
\]

\begin{enumerate}
\def\labelenumi{\alph{enumi})}
\setcounter{enumi}{1}
\tightlist
\item
  Let \(((\mathbf{P}_i)_{i\in I}, (\varphi_{i,j})_{i < j})\) be an inductive system of projective A-modules such that for every pair \((i,j)\) with \(i < j\), \(\varphi_{i,j}\) is injective and the module \(\operatorname{Coker}(\varphi_{i,j})\) is projective. Show that the module \(\mathbf{P} = \varinjlim_{I} \mathbf{P}_i\) is projective. (Let \(0 \to \mathbf{M}' \to \mathbf{M} \to \mathbf{M}'' \to 0\) be an exact sequence of A-modules. Then
\end{enumerate}

\[
0 \to \operatorname{Hom} _ {\mathbf {A}} (\mathrm{P} _ {i}, \mathrm{M} ^ {\prime}) \to \operatorname{Hom} _ {\mathbf {A}} (\mathrm{P} _ {i}, \mathrm{M}) \to \operatorname{Hom} _ {\mathbf {A}} (\mathrm{P} _ {i}, \mathrm{M} ^ {\prime \prime}) \to 0
\]

is an exact sequence of projective systems of commutative groups and the projective system \(i \mapsto \mathrm{Hom}_{\mathbf{A}}(\mathbf{P}_{i}, \mathbf{M}')\) has the Mittag-Leffler property for the discrete topology (TG, II, p.18). From this one will deduce that the sequence of projective limits is exact.)

\begin{enumerate}
\def\labelenumi{\alph{enumi})}
\setcounter{enumi}{2}
\tightlist
\item
  Let \((\mathbf{Q}_i)_{i\in I}\) be an inductive system of projective A-modules. For every \(i\in \mathbf{I}\), let \(\mathbf{R}_i = \bigoplus_{j\leqslant i}\mathbf{Q}_j\) and for \(i < i'\) let \(\psi_{i,i'}: \mathbf{R}_i \to \mathbf{R}_{i'}\) be the canonical injection. Define an exact sequence of inductive systems \(0 \to (\mathbf{P}_i) \to (\mathbf{R}_i) \to (\mathbf{Q}_i) \to 0\) such that \((\mathbf{P}_i)\) satisfies the hypotheses of \(b\). Deduce from this another proof of \(a\).)
\end{enumerate}

\begin{enumerate}
\def\labelenumi{\arabic{enumi})}
\setcounter{enumi}{7}
\tightlist
\item
  Let A be a ring and n an integer ≥ 1. Show that the following conditions are equivalent:
\end{enumerate}

\begin{enumerate}
\def\labelenumi{(\roman{enumi})}
\item
  We have dh(A) ≤ n.
\item
  For every cyclic A-module M, we have \(\mathrm{dp}_{\mathbf{A}}(\mathbf{M}) \leqslant n\) .
\item
  For every ideal I of A, we have \(\mathrm{dp}_{\mathbf{A}}(\mathbf{I}) \leqslant n - 1\) .
\end{enumerate}

(Use A, X, p.138, prop. 4 to prove the equivalence (i) \(\Leftrightarrow\) (ii) and A, X, p.93, prop. 11 to prove (i) \(\Leftrightarrow\) (iii).

\begin{enumerate}
\def\labelenumi{\arabic{enumi})}
\setcounter{enumi}{8}
\tightlist
\item
  Let \(\Gamma\) be a totally ordered commutative group and \(\Gamma^{+}\) the set of positive elements of \(\Gamma\). We denote by (GOD) the following property:
\end{enumerate}

(GOD): Every major subset M ⊂ Γ⁺ (VI, § 3, no. 5, def. 2) has a countable subset that is cofinal for the opposite order.

\begin{enumerate}
\def\labelenumi{\alph{enumi})}
\item
  Show that if \(\Gamma\) is of height 1, property (GOD) is satisfied.
\item
  For every integer n, show that there exists a totally ordered group \(\Gamma\) of height n possessing property (GOD) (take \(\Gamma = \mathbf{Z}^{n}\) equipped with the lexicographic order).
\end{enumerate}

10) Let A be a valuation ring whose ordered group \(\Gamma\) of values has property (GOD). a) Show that the homological dimension of A is at most 2. (Note that every ideal of A is a countable inductive limit of principal ideals, and use the results of Exercises 7 and 8.)

\begin{enumerate}
\def\labelenumi{\alph{enumi})}
\setcounter{enumi}{1}
\item
  Every valuation ring of Krull dimension 1 has homological dimension \(\leqslant 2\). It has homological dimension 2 if and only if it is not Noetherian. (Use exerc. 9, a) and A, X, p.208, exerc. 12.)
\item
  For every integer \(n \geqslant 1\), there exists a ring of Krull dimension equal to n and of homological dimension equal to 2. (Use exerc. 9, b) and A, X, p.208, exerc. 12.)
\end{enumerate}

\begin{enumerate}
\def\labelenumi{\arabic{enumi})}
\setcounter{enumi}{10}
\item
  Every regular Noetherian local ring is factorial. (Use VII, § 4, n° 7, cor. 3.)
\item
  Let A be a ring and let a be a finitely generated ideal of A.
\end{enumerate}

\begin{enumerate}
\def\labelenumi{\alph{enumi})}
\item
  Show that if there exists an ideal b of A such that a = a·b, then there exists \(b \in b\) such that \((1 - b)a = (0)\).
\item
  Show that if \(a^{2} = a\), there exists an element \(a \in a\) such that \(a^{2} = a\) and a = Aa.
\item
  We assume that a is maximal and that the A-module \(a/a^{2}\) can be generated by n elements. Show that the ideal a can be generated by \(n + 1\) elements. (One takes elements \(x_{1}, \ldots, x_{n}\) in a whose images in \(a/a^{2}\) generate this latter. Denoting by x the ideal of A generated by \(\{x_{1}, \ldots, x_{n}\}\) ; one will note that in A/x one has \((\mathfrak{a}/\mathfrak{x})^{2} = (\mathfrak{a}/\mathfrak{x})\) .)
\end{enumerate}

\begin{enumerate}
\def\labelenumi{\arabic{enumi})}
\setcounter{enumi}{12}
\tightlist
\item
  Let \(p\) be a prime number and \(f\) an integer \(> 0\). Put \(L = \mathbf{Q}(\zeta)\), where \(\zeta\) is a primitive \(p^f\)-th root of 1, and let \(B = \mathbf{Z}[\zeta]\) be the subring of L generated by \(\zeta\). Show that B is the integral closure of Z in L. (By a discriminant computation, one shows that \(B\left[\frac{1}{p}\right]\) is the integral closure of \(\mathbf{Z}\left[\frac{1}{p}\right]\) (V, § 1, no. 6, lemma 3); denoting \(\overline{\mathbf{B}}\) the integral closure of \(\mathbf{B}\),
\end{enumerate}

we deduce from this that \(\overline{\mathbf{B}}\subset\bigcap_{q}\mathbf{Z}_{q}[\zeta]\) where \(q\) runs through the prime numbers, whence \(\overline{\mathbf{B}}\subset\mathbf{Z}[\zeta].\)

\begin{enumerate}
\def\labelenumi{\arabic{enumi})}
\setcounter{enumi}{13}
\tightlist
\item
  Let \(\rho: A \to B\) be a homomorphism of Noetherian rings making \(B\) an \(A\)-module that is faithfully flat.
\end{enumerate}

\begin{enumerate}
\def\labelenumi{\alph{enumi})}
\item
  If B is regular, then A is regular (cf. p. 94, exercises 2 and 3).
\item
  If A is regular and if, for every maximal ideal m of A, B/mB is regular, then B is regular. (By localizing, reduce to the case where A and B are local and ρ is local. One then has dim(B) = dim(A) + dim(B/mB) (§ 3, no. 4, cor. 1 to prop. 7). Construct a regular sequence for B having dim(B) terms.)
\end{enumerate}

\begin{enumerate}
\def\labelenumi{\arabic{enumi})}
\setcounter{enumi}{14}
\item
  Let X be a Noetherian space and U a subset of X. If for every irreducible closed subset Y of X meeting U, U ∩ Y contains a nonempty open subset of Y, then U is open in X. (Assuming for contradiction that U is not open, consider a minimal closed set Z ⊂ X such that Z ∩ U is not open.)
\item
  Let A be a Noetherian ring. The singular locus of Spec(A), denoted Sing(A), is the set of \(p \in \text{Spec}(A)\) such that \(A_p\) is not regular. We set
\end{enumerate}

\[
\operatorname{Reg} (\mathrm{A}) = \operatorname{Spec} (\mathrm{A}) - \operatorname{Sing} (\mathrm{A}).
\]

Show that the following conditions are equivalent:

\begin{enumerate}
\def\labelenumi{(\roman{enumi})}
\item
  For every quotient ring B of A, Reg(B) is open in Spec(B).
\item
  For every integral quotient ring B of A, Reg(B) is a nonempty open subset of Spec(B).
\item
  For every integral quotient ring B of A, Reg(B) contains a nonempty open subset of Spec(B).
\end{enumerate}

(To prove (i) \(\Rightarrow\) (ii), note that one has \((0)\in\mathrm{Reg}(\mathbf{B})\) . To prove (iii) \(\Rightarrow\) (i), show that for every \(\mathfrak{p}\in\mathrm{Spec}(\mathbf{A})\) such that \(\mathrm{V}(\mathfrak{p})\cap\mathrm{Reg}(\mathbf{A})\neq\emptyset\) , the ring \(\mathbf{A}_{\mathfrak{p}}\) is regular (p.94, exerc. 5). Using (iii) and constructing completely secant sequences, show that \(\mathrm{V}(\mathfrak{p})\cap\mathrm{Reg}(\mathbf{A})\) contains a nonempty open subset of \(\mathrm{V}(\mathfrak{p})\) . Conclude using exerc. 15.)

\begin{enumerate}
\def\labelenumi{\arabic{enumi})}
\setcounter{enumi}{16}
\tightlist
\item
  Let \(\rho: A \to B\) be an injective homomorphism of integral Noetherian rings making \(B\) an \(A\)-module of finite type.
\end{enumerate}

\begin{enumerate}
\def\labelenumi{\alph{enumi})}
\item
  We assume that \(\operatorname{Reg}(\mathbf{B})\) contains a nonempty open subset of \(\operatorname{Spec}(\mathbf{B})\) . Show that \(\operatorname{Reg}(\mathbf{A})\) contains a nonempty open subset of \(\operatorname{Spec}(\mathbf{A})\) (Reduce to the case where \(\mathbf{B}\) is flat over \(\mathbf{A}\) using II, § 5, no. 1, corollary to prop. 2, then in the case where \(\mathbf{B}\) is regular by localizing with respect to a suitable element \(x \in \mathbf{A}\) Conclude using Exerc. 14, a).
\item
  We assume that \(\operatorname{Reg}(\mathbf{A})\) contains a nonempty open subset of \(\operatorname{Spec}(\mathbf{A})\) and that the field \(\mathbf{K}'\) of fractions of \(\mathbf{B}\) is a separable extension of the field \(\mathbf{K}\) of fractions of \(\mathbf{A}\) . Show that \(\operatorname{Reg}(\mathbf{B})\) contains a nonempty open subset of \(\operatorname{Spec}(\mathbf{B})\) . (Let \(\mathbf{Z} \in \mathbf{K}'\) such that \(\mathbf{K}' = \mathbf{K}(\mathbf{Z})\) (A, V, p.39), and let \(\mathrm{P}(\mathrm{X}) \in \mathrm{K}[\mathrm{X}]\) be the minimal polynomial of Z. Replacing A by the ring \(A_{f}\) for a suitable element f of A, show that one can reduce to proving the desired property of Spec(B) in the following case:
\end{enumerate}

\(\alpha)\) \(\mathbf{Z}\in \mathbf{B}\) and A is regular;

\(\beta)\) we have \(\mathrm{P}(\mathrm{X})\in \mathrm{A}[\mathrm{X}]\) ;

\(\gamma)\) we have \(\mathbf{B} = \mathbf{A}[\mathbf{X}] / (\mathbf{P}(\mathbf{X}))\) (II, § 5, no. 1, prop. 2);

δ) P'(Z) is invertible in B (observe that P'(Z) ≠ 0 in K' because K' is separable over K. Then use II, § 5, no. 1, corollary to prop. 3).

Show that in this case, for every maximal ideal m of A, B/mB is an étale (A/m)-algebra (A, V, p. 32), hence regular. Conclude using exerc. 14, b).

18) Let \(k\) be a field and A a \(k\)-algebra of finite type. Show that \(\operatorname{Reg}(A)\) is an open subset of \(\operatorname{Spec}(A)\). (Using exerc. 16, reduce to proving that \(\operatorname{Reg}(A)\) contains a nonempty open subset when A is integral. Using the normalization lemma (V, § 3, no. 1, th. 1) and exerc. 17, a), reduce to proving this last property when A is the integral closure of \(k[X_1, ..., X_n]\) in a finite quasi-Galois extension K of \(k(X_1, ..., X_n)\). Let K' be the purely inseparable closure of \(k(X_1, ..., X_n)\) in K and A' the integral closure of \(k[X_1, ..., X_n]\) in K'. Reduce to proving the desired property for A' by means of exerc. 17, b). Drawing inspiration from the proof of th. 2 of V, § 3, no. 2, embed K' in an extension \(k'(X_1^{q^{-1}}, ..., X_n^{q^{-1}}) = K''\), where \(k'\) is a purely inseparable extension of \(k\), and reduce to the case where A′ is the integral closure of \(k[X_1, ..., X_n]\) in K'' by virtue of exerc. 17, a). Then observe that A' = \(k'[X_1^{q^{-1}}, ..., X_n^{q^{-1}}]\) (V, § 1, No. 3, cor. 2 to prop. 13) and that A' is regular (p. 94, exerc. 6).

\begin{enumerate}
\def\labelenumi{\arabic{enumi})}
\setcounter{enumi}{18}
\item
  Let A be an integral domain of finite type over a field k. The set of maximal ideals \(m \in \text{Specmax}(A)\) such that \(A_{m}\) is regular is open and dense in \(\text{Specmax}(A)\).
\item
  Let \(k\) be a field, \(k_{i}\) (\(i = 1, 2\)) two extensions of \(k\), and \(n_{i}\) (\(i = 1, 2\)) the transcendence degree of \(k_{i}\) over \(k\) if the latter is finite, or the symbol +\(\infty\) otherwise. Show that we have
\end{enumerate}

\[
\dim (k _ {1} \otimes_ {k} k _ {2}) = \inf (n _ {1}, n _ {2})
\]

and that \(k_{1} \otimes_{k} k_{2}\) is Noetherian if one of the extensions \(k_{i}\) is of finite type.

\begin{enumerate}
\def\labelenumi{\arabic{enumi})}
\setcounter{enumi}{20}
\item
  Let I be a directed ordered set, \((\mathbf{A}_{i}, \varphi_{ij})\) an inductive system of rings with inductive limit A. Suppose that, for every i, the ring \(A_{i}\) is Noetherian, that for every pair \((i, j)\) such that i \textless{} j, \(\varphi_{ij}\) makes \(A_{j}\) a flat \(A_{i}\)-module, and that A is Noetherian. Show that if the \(A_{i}\) are regular for all \(i \in I\), then A is regular. (One may first reduce to the case where the \(A_{i}\) and A are local and the \(\varphi_{ij}\) are local. One will then use exercises 3, p.94 and 14, p.96.)
\item
  \begin{enumerate}
  \def\labelenumii{\alph{enumii})}
  \tightlist
  \item
    Let \(k\) be a field, A a \(k\)-algebra of finite type, and \(k'\) a finite separable extension of \(k\). Show that if A is regular, then \(A \otimes_k k'\) is regular. (One may use exercise 14, b), p.96.)
  \end{enumerate}
\end{enumerate}

\begin{enumerate}
\def\labelenumi{\alph{enumi})}
\setcounter{enumi}{1}
\tightlist
\item
  With the same hypotheses on A, let \(k'\) be a finite purely inseparable extension of \(k\). Show that in general \(A \otimes_k k'\) is not regular (take \(A = k'\)).
\end{enumerate}

\begin{enumerate}
\def\labelenumi{\arabic{enumi})}
\setcounter{enumi}{22}
\tightlist
\item
  Let \(k\) be a field of characteristic exponent \(p\), A a \(k\)-algebra of finite type, and \(\mathfrak{p}\) a prime ideal of A. Show that the following conditions are equivalent:
\end{enumerate}

\begin{enumerate}
\def\labelenumi{(\roman{enumi})}
\item
  the ring \(A_{\mathfrak{p}} \otimes_{k} k'\) is regular for every extension \(k'\) of k;
\item
  the ring \(A_{\mathfrak{p}} \otimes_{k} k^{p^{-\infty}}\) is regular;
\item
  the ring \(\mathbf{A}_{\mathfrak{p}} \otimes_{k} k'\) is regular for every finite and purely inseparable extension \(k'\) of \(k\) .
\end{enumerate}

(To prove (ii) \(\Rightarrow\) (iii), use exercise 14, a), p.96. To prove (iii) \(\Rightarrow\) (i), it suffices to show that \(A_{\mathfrak p} \otimes_k k'\) is regular when \(k'\) is a finitely generated extension of \(k\) (Exerc. 21), and when \(k'\) is a purely inseparable extension of a purely transcendental extension (Exerc. 22 and 14, a), and is bounded by a finite quasi-Galois extension of a purely transcendental extension. Then bound \(k'\) by an extension \(k''(\mathrm{X}_1, ..., \mathrm{X}_n)\), where \(k''\) is a finite purely inseparable extension of \(k\), and reduce to the case \(k' = k''(\mathrm{X}_1, ..., \mathrm{X}_n)\). By exercise 14, a), setting \(\mathbf{B} = \mathbf{A}_{\mathfrak{p}} \otimes_k k''\), we obtain a regular algebra by (iii), and \(\mathbf{A}_{\mathfrak{p}} \otimes_k k' = \mathbf{B} \otimes_{k''} k''(\mathbf{X}_1, ..., \mathbf{X}_n)\) is a regular algebra by exerc. 6, e), p.94.)

We say that \(p \in \text{Spec}(A)\) is smooth if it possesses the equivalent properties above.

\begin{enumerate}
\def\labelenumi{\arabic{enumi})}
\setcounter{enumi}{23}
\item
  Let \(k\) be a field, A a \(k\)-algebra, \(k'\) a purely inseparable extension of \(k\), and \(\mathbf{A}' = \mathbf{A} \otimes_k k'\). Show that the canonical map \(\operatorname{Spec}(\mathbf{A}') \to \operatorname{Spec}(\mathbf{A})\) is a homeomorphism. (When A is a field, \(\mathbf{A}'\) has dimension 0 (exerc. 20); one will show that it possesses only one prime ideal. From this one will deduce that for every prime ideal \(\mathfrak{p}\) of A, the radical \(\mathfrak{p}'\) of \(\mathfrak{p}\mathbf{A}'\) is a prime ideal and that the map \(\mathfrak{p} \mapsto \mathfrak{p}'\) is continuous.)
\item
  Let \(k\) be a field, A a \(k\)-algebra of finite type, and \(\operatorname{Lis}(A)\) the set of \(\mathfrak{p} \in \operatorname{Spec}(A)\) which are smooth (exerc. 23). Show that \(\operatorname{Lis}(A)\) is open in \(\operatorname{Spec}(A)\). (Use exercs. 18, 23 and 24.)
\item
  Let \(k\) be a field and A a \(k\)-algebra that is an integral domain of finite type. Show that Lis(A) is dense if and only if A is a separable \(k\)-algebra; equivalently, the field of fractions of A is a separable extension of \(k\) (A, V, p.115).
\item
  Let \(k\) be a field and A a \(k\)-algebra of finite type that is integral and separable. The set of maximal ideals of A that are smooth is an open dense subset of Specmax(A).
\item
  Let \(k\) be a field and \(k'\) an extension of \(k\). Show that for every extension \(k''\) of \(k\), \(k' \otimes_k k''\) is a regular Noetherian ring if and only if \(k'\) is a separable extension of finite type of \(k\).
\end{enumerate}

29) Let k be a field, A a complete Noetherian local k-algebra, \(\kappa_{A}\) the residue field of A.

\begin{enumerate}
\def\labelenumi{\alph{enumi})}
\item
  Suppose that \(\kappa_{\mathbf{A}}\) is a separable extension of \(k\). Show that there exists a homomorphism of \(k\)-algebras from \(\kappa_{\mathbf{A}}\) to A, which is a section of the canonical homomorphism from A to \(\kappa_{\mathbf{A}}\) (IX, § 3, n° 2, prop. 1).
\item
  Suppose that \(\kappa_{\mathbf{A}}\) is a purely inseparable extension of \(k\). Show that there does not necessarily exist a section of \(k\)-algebras from \(\kappa_{\mathbf{A}}\) into A. (Take a field \(k\) of characteristic 2, an element \(a\) of \(k\) which is not a square (for example \(k = \mathbf{F}_2(\mathrm{T})\), \(a = \mathrm{T}\)), the ideal \(\mathfrak{p}\) of \(k[\mathrm{X}]\) generated by \(\mathrm{X}^2 + a\), and set \(\mathrm{A} = \widehat{k[\mathrm{X}]_{\mathfrak{p}}}\).)
\item
  We assume A is regular and \(\kappa_{\mathbf{A}}\) is a separable extension of \(k\). Show that A is isomorphic as a \(k\)-algebra to \(\kappa_{\mathbf{A}}[[\mathrm{T}_{1}, ..., \mathrm{T}_{n}]]\), where \(n = \dim(\mathbf{A})\).
\item
  Show that this is not necessarily the case when \(\kappa_{\mathrm{A}}\) is not a separable extension of \(k\) .
\end{enumerate}

\begin{enumerate}
\def\labelenumi{\arabic{enumi})}
\setcounter{enumi}{29}
\item
  Let \(k\) be a field and A a \(k\)-algebra, local, Noetherian, and complete. One says that A is formally smooth if, for every finite extension \(k'\) of \(k\), the ring \(A \otimes_k k'\) is regular. Show that A is formally smooth if and only if, for every finite purely inseparable extension \(k'\) of \(k\), the ring \(A \otimes_k k'\) is regular (one may use exerc. 14, p.96 and imitate the proof of exerc. 23, p.97).
\item
  Let \(k\) be a field.
\end{enumerate}

\begin{enumerate}
\def\labelenumi{\alph{enumi})}
\item
  Let A be a \(k\)-algebra, Noetherian, local, complete, and regular, whose residue field \(\kappa_{\mathrm{A}}\) is a separable extension of \(k\). Show that A is formally smooth and isomorphic to \(\kappa_{\mathrm{A}}[[\mathbf{T}_{1}, ..., \mathbf{T}_{n}]]\).
\item
  Let B be a \(k\)-algebra of finite type, \(\mathfrak{p} \in \operatorname{Spec}(\mathbf{B})\) a smooth point (p.97, exerc. 23). Show that \(\hat{\mathbf{B}}_{\mathfrak{p}}\) is a formally smooth \(k\)-algebra.
\item
  Let \(\kappa\) be a finitely generated extension of \(k\) and let \(n\in\mathbf{N}\). Show that there exists a formally smooth \(k\)-algebra with residue field \(\kappa\) and of dimension \(n\).
\item
  Let A be a formally smooth k-algebra whose residue field \(\kappa_{A}\) is not a separable extension of k. Show that A is not isomorphic as a k-algebra to \(\kappa_{A}[[T_{1},...,T_{n}]]\).
\end{enumerate}

32) Let \(k\) be a field, K a finitely generated extension of \(k\), \(n\) a positive integer, and A and B two formally smooth \(k\)-algebras of dimension \(n\), such that the extensions \(\kappa_{A}\) and \(\kappa_{B}\) of \(k\) are isomorphic to K. It is proposed to show that A and B are isomorphic \(k\)-algebras.

\begin{enumerate}
\def\labelenumi{\alph{enumi})}
\item
  Reduce to the case where K is a finite purely inseparable extension of k (exercise 29, a)).
\item
  Consider the completed algebra \(C = A \hat{\otimes}_{k} B\) of \(A \otimes_{k} B\) with respect to the ideal
\end{enumerate}

\[
\left(\mathbf {A} \otimes_ {k} \mathfrak {m} _ {\mathbf {B}}\right) + \left(\mathfrak {m} _ {\mathbf {A}} \otimes_ {k} \mathbf {B}\right).
\]

Show that C is a Noetherian, local, complete, and regular \(k\)-algebra. (For this last point, one uses the canonical homomorphism \(\rho\) from A to C. We will show that it makes C a flat A-module and that \(C \otimes_{\mathbf{A}} \kappa_{\mathbf{A}}\) is isomorphic to \(\kappa_{\mathbf{A}} \otimes_{k} \mathbf{B}\). One will then use exerc. 14, p.96.)

\begin{enumerate}
\def\labelenumi{\alph{enumi})}
\setcounter{enumi}{2}
\tightlist
\item
  Show that \(\rho\) induces an isomorphism on the residue fields and an injection of \(m_{\mathbf{A}}/m_{\mathbf{A}}^{2}\) into \(m_{\mathbf{C}}/m_{\mathbf{C}}^{2}\). Then construct a retraction \(\pi\) of \(\rho\) which is a homomorphism of \(k\)-algebras. Show that \(\pi \circ \rho': B \to A\) is an isomorphism of \(k\)-algebras. (We have denoted by \(\rho'\) the canonical homomorphism from B to C.)
\end{enumerate}

\exercisesection{DIMENSION OF GRADED RINGS}

\begin{enumerate}
\def\labelenumi{\arabic{enumi})}
\tightlist
\item
  Let A be a ring, and let p be a graded prime ideal of A{[}T{]}, and \(p_{0}\) the prime ideal \(p \cap A\) of A.
\end{enumerate}

\begin{enumerate}
\def\labelenumi{\alph{enumi})}
\item
  If \(T \in p\) , then \(p = p_{0} + T \cdot A[T]\) .
\item
  If T \(\notin\) p, then \(\mathfrak{p} = \mathfrak{p}_0\) . A{[}T{]}.
\item
  Deduce that, in the case \(a)\) we have \(\operatorname{htgr}(\mathfrak{p}) = \operatorname{ht}(\mathfrak{p}_0) + 1\) , and in the case \(b)\) , one has \(\operatorname{htgr}(\mathfrak{p}) = \operatorname{ht}(\mathfrak{p}_0)\) .
\item
  Show that \(\dim \operatorname{gr}(\mathbf{A}[\mathbf{T}]) = \dim (\mathbf{A}) + 1\) .
\item
  Deduce from this an example of a graded ring B such that dim(B) ≠ dimgr(B) (p.84, exerc. 7).
\end{enumerate}

\begin{enumerate}
\def\labelenumi{\arabic{enumi})}
\setcounter{enumi}{1}
\tightlist
\item
  Let A be a ring, and \(\mathcal{F} = (\mathcal{F}_i)_{i \in \mathbf{Z}}\) a decreasing filtration of A such that \(\mathcal{F}_i = \mathbf{A}\) for \(i \leqslant 0\). Let \(\mathcal{A}\) be the subring \(\bigoplus_{i \in \mathbf{Z}} \mathcal{F}_i \cdot v^{-i}\) of \(\mathbf{A}[v, v^{-1}]\), and let \(\mathrm{gr}_{\mathcal{F}}(\mathbf{A})\) be the associated graded ring
\end{enumerate}

of the filtered ring (A, F).

\begin{enumerate}
\def\labelenumi{\alph{enumi})}
\item
  Show that the homomorphism \(f: \mathcal{A} \to \mathrm{gr}_{\mathcal{F}}(\mathbf{A})\) defined by \(f\left(\sum_{i} a_i v^{-i}\right) = \sum_{i} \overline{a}_i\), where \(\overline{a}_i\) is the class of \(a_i\) in \(\mathcal{F}_i / \mathcal{F}_{i+1}\), induces an isomorphism of \(\mathcal{A}/v. \mathcal{A}\) with \(\mathrm{gr}_{\mathcal{F}}(\mathbf{A})\).
\item
  Show that for every invertible element \(v_{0}\) of A, the homomorphism \(e_{v_{0}}: \mathcal{A} \to \mathbf{A}\) defined by
\end{enumerate}

\[
e _ {v _ {0}} (\sum a _ {i} v ^ {- i}) = \sum a _ {i} v _ {0} ^ {- i}
\]

induces an isomorphism of \(\mathcal{A}/(v - v_{0})\mathcal{A}\) with A.

\begin{enumerate}
\def\labelenumi{\alph{enumi})}
\setcounter{enumi}{2}
\item
  Show that \(\mathcal{A}/\mathrm{A}[v]\) is a torsion \(\mathrm{A}[v]\)-module.
\item
  Now assuming that A contains a field k, show that the composite morphism
\end{enumerate}

\[
k [ v ] \rightarrow \mathrm{A} [ v ] \rightarrow \mathcal {A} \quad (\text {where } \mathrm{A} [ v ] = \bigoplus_ {i \leqslant 0} \mathcal {F} _ {i}. v ^ {- i})
\]

makes \(\mathcal{A}\) into a k{[}v{]}-torsion-free module. Deduce that \(\mathcal{A}\) is a flat k{[}v{]}-module.

\begin{enumerate}
\def\labelenumi{\alph{enumi})}
\setcounter{enumi}{4}
\tightlist
\item
  Suppose that A is local, with maximal ideal m, and contains a field k such that A is the direct sum of k and m (as an additive group). Show that there exists an ideal S of A such that \(\mathcal{A}\) is the direct sum of k{[}v{]} and S (we have set \(\mathcal{F}_{i} = m^{i}\) for every \(i \in \mathbf{Z}\)).
\end{enumerate}

\begin{enumerate}
\def\labelenumi{\arabic{enumi})}
\setcounter{enumi}{2}
\tightlist
\item
  Let us resume the notation of the previous exercise. Let M be an A-module, equipped with a filtration \((\mathbf{K}_i)_{i\in\mathbf{Z}}\) such that \(\mathbf{K}_i = \mathbf{M}\) for \(i \leqslant 0\).
\end{enumerate}

\begin{enumerate}
\def\labelenumi{\alph{enumi})}
\tightlist
\item
  Show that the graded A-module \(\mathcal{M} = \sum_{i\in \mathbf{Z}}\mathrm{K}_{i}v^{-i}\) has a natural structure of \(\mathcal{A}\) -graded module.
\end{enumerate}



\begin{enumerate}
\def\labelenumi{\alph{enumi})}
\setcounter{enumi}{1}
\item
  Show that the graded A-algebra \(\mathcal{A}\) is generated by its elements of degree 1 and -1 if and only if there exists an ideal q of A such that \(\mathcal{F}_i = \mathfrak{q}^i\) ( \(i \in \mathbf{Z}\) ).
\item
  Assuming this last condition is satisfied, show that the filtration \((\mathbf{K}_i)_{i\in \mathbb{Z}}\) is q-good if and only if the \(\mathcal{A}\) -graded module \(\mathcal{M}\) is of finite type.
\item
  Show that if the filtration \((\mathbf{K}_{i})_{i\in\mathbf{Z}}\) is q-good, the A-module M is v-torsion-free, and conversely, given a finitely generated A-module M without v-torsion, one can equip it with a q-good filtration such that the associated A-module is M.
\end{enumerate}

\begin{enumerate}
\def\labelenumi{\arabic{enumi})}
\setcounter{enumi}{3}
\tightlist
\item
  Let us take up the notations and hypotheses of exercise 2. We say that the filtration \(\mathcal{F}\) is a filtration by quasi-powers if there exists an integer \(k > 0\) such that
\end{enumerate}

\[
\mathcal {F} _ {n} = \sum_ {i = 1} ^ {k} \mathcal {F} _ {i}. \mathcal {F} _ {n - i} \quad \text { for } \quad n \geqslant 1.
\]

Show that if the ring A is Noetherian, the following conditions are equivalent:

\begin{enumerate}
\def\labelenumi{(\roman{enumi})}
\item
  The filtration \(\mathcal{F}\) is a filtration by quasi-powers.
\item
  The ring \(\mathcal{A}\) is Noetherian.
\item
  The ideal \(\mathcal{N} = \sum_{i > 1} v^{-i}\mathcal{F}_i\) of \(\mathcal{A}\) is of finite type.
\item
  There exists an integer \(k \geqslant 1\) such that \(\mathcal{F}\) is the smallest filtration of A for the order relation “\(\mathcal{H}_n \subset \mathfrak{G}_n\) for every \(n\)”, whose first terms are \(\mathcal{F}_1, ..., \mathcal{F}_k\).
\item
  The A-algebra \(\mathcal{A}\) is of finite type.
\end{enumerate}

\begin{enumerate}
\def\labelenumi{\arabic{enumi})}
\setcounter{enumi}{4}
\tightlist
\item
  Let \(\rho: A \to B\) be an injective homomorphism of Noetherian rings making \(B\) an \(A\)-algebra of finite type. We call the transcendence degree of \(B\) over \(A\), and denote it by \(d_A(B)\), the largest integer \(d\) such that there exists an injective homomorphism of \(A\)-algebras \(b: A[X_1, \ldots, X_d] \to B\). a) Show that if \(\sigma: B \to C\) satisfies the same hypotheses as \(\rho\), one has the inequality
\end{enumerate}

\[
\mathrm{d} _ {\mathbf {A}} (\mathbf {C}) \geqslant \mathrm{d} _ {\mathbf {A}} (\mathbf {B}) + \mathrm{d} _ {\mathbf {B}} (\mathbf {C}).
\]

\begin{enumerate}
\def\labelenumi{\alph{enumi})}
\setcounter{enumi}{1}
\item
  Show that if A and B are integral domains, \(d_{\mathbf{A}}(\mathbf{B})\) is equal to the transcendence degree of the field of fractions of B over that of A. Moreover, if in a) one assumes C is an integral domain, one has the equality \(d_{\mathbf{A}}(\mathbf{C}) = d_{\mathbf{A}}(\mathbf{B}) + d_{\mathbf{B}}(\mathbf{C})\).
\item
  Show that if B is a subalgebra of A{[}v, \(v^{-1}\) {]} containing A{[}v{]}, where v is an indeterminate, we have \(d_{\mathbf{A}}(\mathbf{B}) = 1\) .
\end{enumerate}

6) Let \(\rho: A \to B\) be an injective homomorphism of Noetherian rings, making \(B\) an \(A\)-algebra of finite type. Let \(q\) be a prime ideal of \(B\) distinct from \(B\), and set \(p = \rho^{-1}(q)\). a) Show that, with the notations of the preceding exercise, one has the inequality

\[
(*) \quad \mathrm{ht} (\mathfrak {q}) + \mathrm{d} _ {\mathrm{A} / \mathfrak {p}} (\mathrm{B} / \mathfrak {q}) \leqslant \mathrm{ht} (\mathfrak {p}) + \mathrm{d} _ {\mathrm{A}} (\mathrm{B}).
\]

(One may proceed by induction on \(d_{A}(B)\) .)

\begin{enumerate}
\def\labelenumi{\alph{enumi})}
\setcounter{enumi}{1}
\item
  Show that if the A-algebra B is generated by d elements and \(\mathrm{d}_{\mathbf{A}}(\mathbf{B}) = d\) one has equality in (*).
\item
  Show that if A is an integral domain and the polynomial ring A{[}T \(_{1}\) , \ldots, T \(_{n}\) {]} is catenary for every n, equality holds in (*) for every integral A-algebra B of finite type and every prime ideal q of B distinct from B.
\end{enumerate}

\begin{enumerate}
\def\labelenumi{\arabic{enumi})}
\setcounter{enumi}{6}
\tightlist
\item
  Let A be a Noetherian ring and \(\mathcal{F} = (\mathcal{F}_i)_{i\in \mathbf{Z}}\) a decreasing filtration of A such that \(\mathcal{F}_0 = \mathbf{A}\), \(\mathcal{F}_1 \neq \mathbf{A}\), and the A-algebra \(\mathbf{B} = \sum_{i=1}^{\infty} \mathcal{F}_i v^{-i}\) is of finite type.
\end{enumerate}

We propose to show that we have \(\dim(B) = \dim(A) + 1\) and to use this fact to generalize the result of § 6, no. 3, corollary of prop. 5. One will use the results of the two preceding exercises.

\begin{enumerate}
\def\labelenumi{\alph{enumi})}
\tightlist
\item
  Let M be a maximal ideal of B. Show that we have
\end{enumerate}

\[
\operatorname{ht} (\mathfrak {M}) \leqslant \dim (\mathrm{A}) + 1,
\]

and deduce from it the inequality

\[
\dim (\mathbf {B}) \leqslant \dim (\mathbf {A}) + 1.
\]

\begin{enumerate}
\def\labelenumi{\alph{enumi})}
\setcounter{enumi}{1}
\tightlist
\item
  Using the fact that every element of the A{[}v{]}-module A{[}v, \(v^{-1}\) {]}/A{[}v{]} is annihilated by a power of v, show that we have
\end{enumerate}

\[
\dim (\mathrm{A} [ v, v ^ {- 1} ]) = \dim (\mathrm{A} [ v ]) = \dim (\mathrm{A}) + 1,
\]

and deduce that we have \(\dim(\mathbf{B}) \geqslant \dim(\mathbf{A}) + 1\) , whence finally \(\dim(\mathbf{B}) = \dim(\mathbf{A}) + 1\) (one may use the fact that the ring \(\mathbf{A}[v, v^{-1}]\) is a ring of fractions of \(\mathbf{B}\) ). c) Show that the minimal prime ideals of \(\mathbf{B}\) are the ideals of the form:

\[
\sum_ {i \in \mathbf {Z}} \left(\mathfrak {p} \cap \mathcal {F} _ {i}\right) v ^ {- i}
\]

where p is a minimal prime ideal of A.

\begin{enumerate}
\def\labelenumi{\alph{enumi})}
\setcounter{enumi}{3}
\item
  Show that one has \(v\mathbf{B} \neq \mathbf{B}\) and that \(v\) does not belong to any of the minimal prime ideals of B. Deduce the inequality \(\dim(\mathbf{B}/v\mathbf{B}) \leqslant \dim(\mathbf{A})\) (one may use § 1 and § 3, no. 3, cor. 1 of prop. 4).
\item
  Suppose now that there exists a maximal ideal m of A containing \(\mathcal{F}_1\) and such that ht(m) = dim(A). Show that in this case we have dim(B/vB) ≥ dim(A) and therefore dim(B/vB) = dim(A). (One may reduce to the case where A is a local ring with maximal ideal m, show that the ideal M of B generated by v and \(\sum_{i\in\mathbf{Z}}(\mathcal{F}_i\cap m)v^{-i}\) is a maximal ideal of B of height dim(A) + 1, and localize B at M.)
\end{enumerate}





\begin{enumerate}
\def\labelenumi{\alph{enumi})}
\setcounter{enumi}{5}
\tightlist
\item
  If A is a Noetherian local ring, and F is a decreasing filtration of A such that \(F_{0} = A\) , \(F_{1} \neq A\) , we have the equality
\end{enumerate}

\[
\dim (\operatorname{gr} (A)) = \dim (A).
\]

\begin{enumerate}
\def\labelenumi{\arabic{enumi})}
\setcounter{enumi}{7}
\tightlist
\item
  Let A be a Noetherian ring, and let B be a graded sub-A-algebra of the polynomial ring A{[}T{]}.
\end{enumerate}

\begin{enumerate}
\def\labelenumi{\alph{enumi})}
\tightlist
\item
  Show that we have the inequalities
\end{enumerate}

\[
\dim (\mathbf {A}) \leqslant \dim (\mathbf {B}) \leqslant \dim (\mathbf {A}) + 1.
\]

\begin{enumerate}
\def\labelenumi{\alph{enumi})}
\setcounter{enumi}{1}
\tightlist
\item
  Let \(B_{i}\) be the set of \(b \in A\) such that \(bT^{i} \in B\). Show that if we have \(\dim(B) = \dim(A)\), there exists a positive integer k such that for all \(i \geqslant k\), \(B_{i}\) is contained in the intersection of the minimal prime ideals of A.
\end{enumerate}

\begin{enumerate}
\def\labelenumi{\arabic{enumi})}
\setcounter{enumi}{8}
\item
  Let A be a ring and M a finitely presented A-module. For every \(\mathfrak{p} \in \operatorname{Spec}(A)\), we set \(\kappa(\mathfrak{p}) = A_{\mathfrak{p}} / \mathfrak{p} \cdot A_{\mathfrak{p}}\). Show that for all \(k \in \mathbf{N}\), the set of \(\mathfrak{p} \in \operatorname{Spec}(A)\) such that \(\dim_{\kappa(\mathfrak{p})} (M \otimes_A \kappa(\mathfrak{p})) \geqslant k\) is a closed subset of \(\operatorname{Spec}(A)\). (One may take a presentation \(\mathrm{A}^m \xrightarrow{u} \mathrm{A}^n \longrightarrow \mathrm{M} \longrightarrow 0\) of M and consider the minors of the matrix of \(u\).)
\item
  Let \(\rho: A \to B\) be a ring homomorphism. For every \(\mathfrak{p} \in \operatorname{Spec}(B)\), we set
\end{enumerate}

\[
\mathrm{d} (\mathfrak {p}) = \dim_ {\mathfrak {p}} \left(\mathbf {B} \otimes_ {\mathbf {A}} \kappa \left(\rho^ {- 1} (\mathfrak {p})\right)\right).
\]

The aim of this exercise is to prove that when B is a finitely generated A-algebra, the map \(p \mapsto d(p)\) from \(\operatorname{Spec}(B)\) to \(\mathbf{R}\) is upper semicontinuous (“Chevalley’s theorem”).

Denote by \(\delta \subset \mathbf{B} \otimes_{\mathbf{A}} \mathbf{B}\) the kernel of the canonical A-homomorphism \(\mathbf{B} \otimes_{\mathbf{A}} \mathbf{B} \to \mathbf{B}\). For every integer \(r\), set

\[
\mathrm{P} _ {\mathbf {B} / \mathbf {A}} ^ {r} = (\mathbf {B} \otimes_ {\mathbf {A}} \mathbf {B}) / \delta^ {r}.
\]

The B ⊗\_A B-algebra P \(_{B/A}^{r}\) is called the algebra of principal parts of order r. We consider B ⊗\_A B as a B-algebra via the homomorphism b ↦ b ⊗ 1, so that P \(_{B/A}^{r}\) is, for every r, a B-algebra and the canonical quotient homomorphisms P \(_{B/A}^{r+1}\) → P \(_{B/A}^{r}\) define a projective system of B-algebras.

\(^{1}\) This is the dimension as a vector space over the field \(\kappa(\mathfrak{p})\).

\begin{enumerate}
\def\labelenumi{\alph{enumi})}
\tightlist
\item
  Let \(p: \mathbf{A} \to \mathbf{A}'\) be a ring homomorphism. Set \(\mathbf{B}' = \mathbf{B} \otimes_{\mathbf{A}} \mathbf{A}'\) and denote by \(q: \mathbf{B} \to \mathbf{B}'\) the canonical homomorphism. Let \(r\) be an integer. Show that the unique homomorphism of \(\mathbf{B}'\)-algebras
\end{enumerate}

\[
\mathbf {P} _ {\mathbf {B} / \mathbf {A}} ^ {r} \otimes_ {\mathbf {B}} \mathbf {B} ^ {\prime} \rightarrow \mathbf {P} _ {\mathbf {B} ^ {\prime} / \mathbf {A} ^ {\prime}} ^ {r}
\]

which associates to the class of \((b_{1} \otimes b_{2}) \otimes 1\) (\(b_{i} \in \mathbf{B}\)) the class of \(q(b_{1}) \otimes q(b_{2})\) is an isomorphism. b) Let S be a multiplicative subset of B. Show that the canonical homomorphism

\[
\mathrm{S} ^ {- 1} \mathrm{P} _ {\mathrm{B/A}} ^ {r} \rightarrow \mathrm{P} _ {\mathrm{S} ^ {- 1} \mathrm{B/A}} ^ {r}
\]

is an isomorphism.

\begin{enumerate}
\def\labelenumi{\alph{enumi})}
\setcounter{enumi}{2}
\item
  Assume that B is a finitely presented A-algebra, that is, isomorphic to a quotient of \(\mathrm{A}[\mathrm{T}_{1},\ldots,\mathrm{T}_{n}]\) by a finitely generated ideal. Show that for every \(r\in\mathbf{N}\) , \(\mathbf{P}_{\mathrm{B}/\mathrm{A}}^{r}\) is a finitely presented B-module. For every \(\mathfrak{p}\in\operatorname{Spec}(\mathbf{B})\) , we set \(\mathrm{P}_{\mathrm{B}/\mathrm{A}}^{\infty}[\mathfrak{p}]=\lim_{r}\mathrm{P}_{\mathrm{B}/\mathrm{A}}^{r}\otimes_{\mathrm{B}}\kappa(\mathfrak{p})\) . Show that \(\mathrm{P}_{\mathrm{B}/\mathrm{A}}^{\infty}[\mathfrak{p}]\) is a \(\kappa(\mathfrak{p})\) -algebra that is local and complete with residue field \(\kappa(\mathfrak{p})\) . Let \(k\in\mathbf{N}\) . Show that the set of \(\mathfrak{p}\in\operatorname{Spec}(\mathbf{B})\) such that \(\dim(\mathrm{P}_{\mathrm{B}/\mathrm{A}}^{\infty}[\mathfrak{p}])\geqslant k\) is a closed subset of \(\operatorname{Spec}(\mathbf{B})\) . (For every \(r\in\mathbf{N}\) , let \(\mathrm{F}_{r,k}\) the set of \(\mathfrak{p}\) such that \(\dim_{\kappa(\mathfrak{p})}(\mathrm{P}_{\mathrm{B}/\mathrm{A}}^{r}\otimes\kappa(\mathfrak{p}))\geqslant\binom{r+k-1}{r-1}\) . We shall show that \(\mathrm{F}_{r,k}\) is closed (Exercise 9) and consequently that \(\Phi_{k}=\bigcap_{r}\mathrm{F}_{r,k}\) is closed. One will then apply § 6, no. 3, Cor. 2 to Prop. 6.)
\item
  Let m be an ideal of B and \(\mu\subset(\mathrm{B/m})\otimes_{\mathrm{A}}\mathrm{B}\) the kernel of the canonical mapping \((\mathrm{B/m})\otimes_{\mathrm{A}}\mathrm{B}\to\mathrm{B/m}\). Show that we have a canonical isomorphism
\end{enumerate}

\[
\mathrm{P} _ {\mathrm{B/A}} ^ {r} \otimes_ {\mathrm{B}} (\mathrm{B/m}) \rightarrow ((\mathrm{B/m}) \otimes_ {\mathrm{A}} \mathrm{B}) / \mu^ {r}.
\]

Assume that A is a field, that B is a finitely generated A-algebra, and that m is a maximal ideal of B. Show that \(P_{B/A}^{\infty}[m]\) is isomorphic to the completed localization of \((B/m)\otimes_{A}B\) at the maximal ideal \(\mu\). Deduce that one has \(\dim(P_{B/A}^{\infty}[m])=\dim_{m}(B)\). (Apply the zero theorem and th. 1 of § 2, no. 3. Note that the extension B/m of A is contained in a finite Galois extension of a finite purely inseparable extension of A.)

\begin{enumerate}
\def\labelenumi{\alph{enumi})}
\setcounter{enumi}{4}
\tightlist
\item
  Assume that B is a finitely generated A-algebra. Let \(p \in \operatorname{Spec}(B)\). Show that we have
\end{enumerate}

\[
\dim (\mathrm{P} _ {\mathrm{B/A}} ^ {\infty} [ \mathfrak {p} ]) = \dim \left(\mathrm{B} \otimes_ {\mathrm{A}} \kappa (\rho^ {- 1} (\mathfrak {p}))\right).
\]

First note that for every \(r \in \mathbf{N}\), one has \(\mathbf{P}_{\mathbf{B}/\mathbf{A}}^{\mathbf{r}} \otimes_{\mathbf{B}} \kappa(\mathfrak{p}) = (\mathbf{P}_{\mathbf{B}/\mathbf{A}}^{\mathbf{r}} \otimes_{\mathbf{B}} \mathbf{B}') \otimes_{\mathbf{B}'} \kappa(\mathfrak{p})\) by setting \(\mathbf{B}' = \mathbf{B} \otimes_{\mathbf{A}} \kappa(\rho^{-1}(\mathfrak{p}))\), and hence, by virtue of \(a\) ), we have \(\mathbf{P}_{\mathbf{B}/\mathbf{A}}^{\infty}[\mathfrak{p}] = \mathbf{P}_{\mathbf{B}' / \kappa(\rho^{-1}(\mathfrak{p}))}^{\infty}[\mathfrak{p}\mathbf{B}']\). One thus reduces to proving, in the case where A is a field, the equality

\[
\dim \left(P _ {B / A} ^ {\infty} [ p ]\right) = \dim_ {p} (B).
\]

Observe then that in the case where p is maximal, this equality follows from d). In the general case, let \(q_{1}, \ldots, q_{l}\) be the minimal prime ideals of B contained in p. Show, using the Nullstellensatz, that there exists a dense open set U of \(V(p)\) such that for every maximal ideal m belonging to U, the minimal prime ideals contained in m are contained in U. Using c), show that one may, by shrinking U, assume moreover that for every \(p' \in U\), one has \(\dim(\mathrm{P}_{\mathrm{B/A}}^{\infty}[\mathfrak{p}]) = \dim(\mathrm{P}_{\mathrm{B/A}}^{\infty}[\mathfrak{p}'])\). Conclude.)

\begin{enumerate}
\def\labelenumi{\alph{enumi})}
\setcounter{enumi}{5}
\tightlist
\item
  Assume that B is a finitely generated A-algebra. Show that the function \(\mathfrak{p} \mapsto \mathrm{d}(\mathfrak{p}) = \dim_{\mathfrak{p}} (\mathrm{B} \otimes_{\mathbf{A}} \kappa(\rho^{-1}(\mathfrak{p})))\) is upper semicontinuous. (When B is of finite presentation, use \(e\) and \(c\). In the general case, B is of the form \(\mathrm{A}[\mathrm{T}_1, ..., \mathrm{T}_n]/\mathfrak{J}\). Let \(\mathfrak{J}_{\alpha}\) be the family of finitely generated ideals of \(\mathrm{A}[\mathrm{T}_1, ..., \mathrm{T}_n]\) contained in \(\mathfrak{J}\), ordered by inclusion, and set \(\mathrm{B}_{\alpha} = \mathrm{A}[\mathrm{T}_1, ..., \mathrm{T}_n]/\mathfrak{J}_{\alpha}\). For every \(\alpha\), one has \(\operatorname{Spec}(\mathbf{B}) \subset \operatorname{Spec}(\mathrm{B}_{\alpha}) \subset \operatorname{Spec}(\mathrm{A}[\mathrm{T}_1, ..., \mathrm{T}_n])\), and \(\operatorname{Spec}(\mathbf{B}) = \bigcap \operatorname{Spec}(\mathrm{B}_{\alpha})\). For every \(\mathfrak{p} \in \operatorname{Spec}(\mathbf{B})\), set
\end{enumerate}

\[
\mathrm{d} _ {\alpha} (\mathfrak {p}) = \dim_ {\mathfrak {p}} \left(\mathrm{B} _ {\alpha} \otimes \kappa \left(\rho^ {- 1} (\mathfrak {p})\right)\right).
\]

Show that \(d_{\alpha}(p) \geqslant d(p)\) and that \(d(p) = \inf_{\alpha} d_{\alpha}(p)\) by virtue of the Noetherian character of \(\text{Spec}(\kappa(\rho^{-1}(p))[T_1, \ldots, T_n])\). Deduce that \(p \mapsto d(p)\) is the infimum of a filtered decreasing family of upper semicontinuous functions.)

\begin{enumerate}
\def\labelenumi{\arabic{enumi})}
\setcounter{enumi}{10}
\tightlist
\item
  Let \(\rho: \mathbf{A} \to \mathbf{B}\) be a ring homomorphism, and let \(^a\rho: \operatorname{Spec}(\mathbf{B}) \to \operatorname{Spec}(\mathbf{A})\) be the corresponding map.
\end{enumerate}

\begin{enumerate}
\def\labelenumi{\alph{enumi})}
\item
  Show that for every \(\mathfrak{q} \in \operatorname{Spec}(\mathbf{A})\), \(\operatorname{Spec}(\mathbf{B} \otimes_{\mathbf{A}} \kappa(\mathfrak{q}))\) is identified with \(^a\rho^{-1}(\mathfrak{q})\).
\item
  Suppose that B is a finitely generated A-algebra. Show that p is an isolated point in its fiber \(^a\rho^{-1}(^a\rho(p))\) if and only if \(\dim_p(B \otimes_A \kappa(^a\rho(p)))\) is zero.
\item
  From exercise 10, deduce that under the hypotheses of \(b\), the set of points of \(\operatorname{Spec}(\mathbf{B})\) which are isolated in their fiber is an open subset of \(\operatorname{Spec}(\mathbf{B})\).
\end{enumerate}

\begin{enumerate}
\def\labelenumi{\arabic{enumi})}
\setcounter{enumi}{11}
\tightlist
\item
  Let H be a finitely generated graded algebra, such that \(H_{0}\) is a field, and \((\xi_{a})_{a\in\mathbf{A}}\) homogeneous elements of H that generate H. Suppose there exists a family of strictly positive integers \((d_{i})_{i\in\mathbf{I}}\) such that \(P_{H}=\prod(1-T^{d_{i}})^{-1}\). Denote by \(\delta_{a}>0\) the degree of \(\xi_{a}\).
\end{enumerate}

\begin{enumerate}
\def\labelenumi{\alph{enumi})}
\item
  Show that there exists an injection \(\varphi: I \to A\) such that for every \(i\), \(d_i\) divides \(\delta_{\varphi(i)}\). (One may use § 4, no. 2, th. 1.)
\item
  Show that one has \(\inf_{a \in A} \delta_a \leqslant \inf_{i \in I} d_i\).
\item
  From this, deduce that if the \(\delta_{a}\) are equal to one another, the graded algebra H is regular.
\item
  We assume that A has two elements \(a\) and \(a'\) such that \(\delta_{a} < \delta_{a'}\) and that \(\delta_{a'}\) is not a multiple of \(\delta_{a}\). Show that the graded algebra H is regular.
\end{enumerate}

\exercisesection{MULTIPLICITIES}

\begin{enumerate}
\def\labelenumi{\arabic{enumi})}
\tightlist
\item
  Let \(\Gamma \subset \mathbf{N}\) be a submonoid of \(\mathbf{N}\) such that \(\mathbf{N} - \Gamma\) is a finite set. Let \((a_1, \ldots, a_r)\), with \(0 < a_1 < \cdots < a_r\), be a minimal system of generators of \(\Gamma\), and let \(k\) be a field. Set \(\mathbf{A} = k[\Gamma]\).
\end{enumerate}

\begin{enumerate}
\def\labelenumi{\alph{enumi})}
\item
  Show that A is isomorphic to the sub-k-algebra of \(k[\mathbf{T}]\) generated by \(\mathbf{T}^{a_1}\), \(\ldots\), and \(\mathbf{T}^{a_r}\).
\item
  Let m be the maximal ideal of A generated by \(T^{a_{1}}\), \(\ldots\), and \(T^{a_{r}}\); let \(A_{m}\) be the local ring of A at m, and set \(q = mA_{m}\). Show that we have \(\dim_{k}q/q^{2} = r\) and \(e_{q}(A_{m}) = a_{1}\).
\end{enumerate}

\begin{enumerate}
\def\labelenumi{\arabic{enumi})}
\setcounter{enumi}{1}
\tightlist
\item
  Let \(k\) be a field, \(m\) an integer \(\geqslant 1\), and I an ideal of the polynomial ring \(k[\mathbf{T}_{1},\ldots,\mathbf{T}_{m}]\) generated by monomials represented by points \(\mathbf{M}_{1},\ldots,\mathbf{M}_{s}\) of the “positive quadrant” \(\mathbf{R}_{+}^{m}\) of \(\mathbf{R}^{m}\). Denote by \(\mathcal{N}\) the convex hull in \(\mathbf{R}_{+}^{m}\) of the union of the sets \(\mathbf{M}_{i}+\mathbf{R}_{+}^{m}\) for \(1\leqslant i\leqslant s\).
\end{enumerate}

\begin{enumerate}
\def\labelenumi{\alph{enumi})}
\tightlist
\item
  Set \(m = (T_1, \ldots, T_m)\), \(A = k[T_1, \ldots, T_m]_m\), and \(q = I.A\). Show that \(A/q\) has finite length if and only if the volume of \(\mathbf{R}_+^m - \mathcal{N}\) is finite.
\item
  Prove the equality \(e_q(A) = m! \operatorname{Vol}(\mathbf{R}_+^m - \mathcal{N})\).
\end{enumerate}

\begin{enumerate}
\def\labelenumi{\arabic{enumi})}
\setcounter{enumi}{2}
\tightlist
\item
  Let \(k\) be a field, and let \(f_{1},...,f_{r}\) be homogeneous elements of the graded ring \(k[X_{1},...,X_{n}]\) forming a complete intersection sequence. One says that the quotient ring
\end{enumerate}

\[
\mathbf {H} = k [ \mathrm{X} _ {1},..., \mathrm{X} _ {n} ] / (f _ {1},..., f _ {r})
\]

is a graded complete intersection ring. Let I be a nonzero graded ideal of the ring \(\mathbf{H} = k[\mathbf{X}_1, \ldots, \mathbf{X}_n] / (f_1, \ldots, f_r)\) endowed with the quotient grading. Show that one has \(\dim(\mathbf{I}) = \dim(\mathbf{H})\) . Deduce that one has either \(\dim(\mathbf{H}/\mathbf{I}) < \dim(\mathbf{H})\) , or \(\dim(\mathbf{H}/\mathbf{I}) = \dim(\mathbf{H})\) and \(c_{\mathbf{M}/\mathbf{I}} < c_{\mathbf{M}}\) ( \(\S 4, \text{n}^{\circ} 2\) ). (One will consider the ideal b formed by the \(s \in \mathbf{H}\) such that \(s.\mathbf{I} = \{0\}\) and we will show that if one has \(\dim(\mathbf{I}) < \dim(\mathbf{H})\) , then b contains a non-zero-divisor element, by examining the position of b relative to the minimal prime ideals of H and to \(\mathbf{H}_{+} = \mathbf{H}_{\geqslant 1}\) .)

\begin{enumerate}
\def\labelenumi{\arabic{enumi})}
\setcounter{enumi}{3}
\tightlist
\item
  Let A be a Noetherian local ring, with maximal ideal m.
\end{enumerate}

\begin{enumerate}
\def\labelenumi{\alph{enumi})}
\item
  Suppose that \(\mathrm{gr}_{\mathfrak{m}}(\mathrm{A})\) is a graded complete intersection ring (cf. Exercise 3). Show that the completion \(\widehat{A}\) of A is a complete intersection, that is, the quotient of a regular local ring by an ideal generated by a completely secant sequence.
\item
  Let \(x \in \mathfrak{m}\), and let \(\xi\) be the image of \(x\) in \(\mathfrak{m}^{\nu}/\mathfrak{m}^{\nu+1}\), where \(\nu\) is such that \(x \in \mathfrak{m}^{\nu}\) and \(x \notin \mathfrak{m}^{\nu+1}\). Show that we have \(e_{\mathfrak{m}}(\mathbf{A}/x\mathbf{A}) = \nu e_{\mathfrak{m}}(\mathbf{A})\) if and only if \(\xi\) is not a zero divisor in \(\mathrm{gr}_{\mathfrak{m}}(\mathbf{A})\), and that in this case one has \(\mathrm{gr}_{\mathfrak{m}}(\mathbf{A}/x\mathbf{A}) = \mathrm{gr}_{\mathfrak{m}}(\mathbf{A})/\xi\mathrm{gr}_{\mathfrak{m}}(\mathbf{A})\) and \(x\) is not a zero divisor in A. (Let I be the homogeneous ideal of \(\mathrm{gr}_{\mathfrak{m}}(\mathbf{A})\) consisting of the elements \(\eta \in \mathrm{gr}_{\mathfrak{m}}(\mathbf{A})\) such that \(\xi\eta = 0\). Consider the ideal \(\mathrm{N}_n = \mathfrak{m}^{\nu+n}:x\) of A and the map from \(\mathrm{I}_n\) into \(\mathrm{N}_{n+1}/\mathfrak{m}^{n+1}\) which sends \(\eta \in \mathrm{I}_n\) to the class modulo \(\mathfrak{m}^{n+1}\) of an element \(y \in A\) lifting \(\eta\). We will show that it is injective, and from this we will deduce the inequality
\end{enumerate}

\[
\operatorname{long} _ {\mathbf {A} / \mathfrak {m}} (\mathbf {I} _ {n}) \leqslant \operatorname{long} _ {\mathbf {A} / \mathfrak {m}} (\mathbf {N} _ {n + 1} / \mathfrak {m} ^ {n + 1});\tag{i}
\]

Moreover, we will prove the equality

\[
(1 - \mathrm{T} ^ {\mathrm{v}}) \mathrm{H} _ {\mathrm{A}} ^ {(1)} = \mathrm{T} ^ {- \mathrm{v}} \mathrm{H} _ {\mathrm{A} / x \mathrm{A}} ^ {(1)} - \sum_ {n \geqslant \mathrm{v}} \operatorname{long} _ {\mathrm{A} / \mathfrak {m}} (\mathrm{N} _ {n} / \mathfrak {m} ^ {n}). \mathrm{T} ^ {n}.\tag{ii}
\]

Finally, we note that \(\xi\) is a zero divisor if and only if I \(\neq\) \{0\}, and using the previous exercise, we deduce from (i) that \(\sum_{n\geqslant 0}\mathrm{long}_{\mathrm{A / m}}(\mathbf{N}_{n + 1} / \mathfrak{m}^{n + 1})\) \(\mathbf{T}^n = \frac{\mathbf{S}(\mathbf{T})}{(1 - \mathbf{T})^\nu}\), with \(\mathbf{S}(1)\neq 0\), and from (ii) the inequality \(e_{\mathfrak{m}}(\mathbf{A} / x\mathbf{A}) > e_{\mathfrak{m}}(\mathbf{A})\). c) Let \((x_{1},\dots,x_{s})\) be a sequence of elements of the maximal ideal m. Suppose we have \(x_{i}\in \mathfrak{m}^{\nu_i}\) and \(x_{i}\notin \mathfrak{m}^{\nu_i + 1}\), and let \(\xi_{i}\) be the class of \(x_{i}\) in \(\mathfrak{m}^{\nu_i} / \mathfrak{m}^{\nu_i + 1}\). Show that we have \(e_{\mathfrak{m}}(\mathbf{A} / x)\geqslant e_{\mathfrak{m}}(\mathbf{A})\nu_{1}\dots \nu_{s}\), and that equality holds if and only if the sequence \((\xi_1,\dots,\xi_s)\) is completely secant in the ring \(\mathrm{gr}_{\mathfrak{m}}(\mathbf{A})\). In this case we have an isomorphism of graded rings from \(\mathrm{gr}_{\mathfrak{m}}(\mathbf{A} / x)\) with \((\mathrm{gr}_{\mathfrak{m}}(\mathbf{A}))/\xi\), where \(\mathbf{x} = \sum_{i}x_{i}\mathbf{A}\) and \(\xi = \sum_{i}\xi_{i}\mathrm{gr}_{\mathfrak{m}}(\mathbf{A})\).

\begin{enumerate}
\def\labelenumi{\arabic{enumi})}
\setcounter{enumi}{4}
\tightlist
\item
  Let \(k\) be a field. Set \(\mathbf{R} = k[[\mathrm{X},\mathrm{Y},\mathrm{Z}]]\) and consider the local ring \(\mathbf{A} = \mathbf{R}/(\mathbf{XY},\mathbf{X}\mathbf{Z})\) with maximal ideal m. Show that one has
\end{enumerate}

\[
\mathrm{H} _ {\mathrm{A}} = (1 - \mathrm{T}) (1 - \mathrm{T} - \mathrm{T} ^ {2}) \mathrm{H} _ {\mathrm{R}} = 1 + \sum_ {n \geqslant 1} (n + 2) \mathrm{T} ^ {n}.
\]

Hence we have \(\dim(A) = 2\), \(\dim_k(m/m^2) = 3\), and \(e_m(A) = 1\). (We have \(H_A = P_G\), where \(G = gr_m(A) = k[X, Y, Z]/(XY, XZ)\). We shall prove the existence of an exact sequence of graded H-modules

\[
0 \rightarrow \mathrm{H} (- 3) \rightarrow \mathrm{H} (- 2) \oplus \mathrm{H} (- 2) \rightarrow \mathrm{H} \rightarrow \mathrm{G} \rightarrow 0
\]

\(\text{where } \mathrm{H} = k[\mathrm{X},\mathrm{Y},\mathrm{Z}].\)

\begin{enumerate}
\def\labelenumi{\arabic{enumi})}
\setcounter{enumi}{5}
\item
  Let \(k\) be a field, A a local, complete, reduced \(k\)-algebra of dimension 1, and \(m\) its maximal ideal. Show that the minimum number of generators of an ideal of A is at most \(e_m(A)\). There exists \(X \in A\) such that the injection of \(k[[X]]\) into A makes A into a \(k[[X]]\)-free module of rank \(e_m(A)\).
\item
  Let A be a local ring with maximal ideal m, and let M be a finitely generated A-module, \((x_{1}, \ldots, x_{d})\) a sequence of elements of m that is secant for the A-module M, and x the ideal \(x_{1}A + \cdots + x_{d}A\).
\end{enumerate}

\begin{enumerate}
\def\labelenumi{\alph{enumi})}
\tightlist
\item
  Suppose the A-module M is faithful. Then the quotient of \(\mathrm{gr}_{x}(\mathbf{A})\) by its nilradical n is isomorphic to \((\mathbf{A}/\mathfrak{m})[\mathbf{X}_{1}, \ldots, \mathbf{X}_{d}]\), and \((\mathrm{gr}_{x}(\mathbf{M}))_{n}\) is a \((\mathrm{gr}_{x}(\mathbf{A}))_{n}\)-module of finite length \(e_{x}(\mathbf{M})\).
\item
  In the general case, \(\mathrm{gr}_{x}(\mathbf{A})\) has a unique minimal prime ideal n, and \((\mathrm{gr}_{x}(\mathbf{M}))_{n}\) is a \(\mathrm{gr}_{x}(\mathbf{A})\)-module of finite length \(e_{x}(\mathbf{M})\).
\end{enumerate}

\begin{enumerate}
\def\labelenumi{\arabic{enumi})}
\setcounter{enumi}{7}
\tightlist
\item
  Let A be a ring and a an ideal of A. Set \(a^{i} = A\) for \(i \leqslant 0\) and consider the A-algebra \(\mathcal{A}(\mathfrak{a}) = \sum_{i \in \mathbf{Z}} \mathfrak{a}^{-i} v^{i} \subset \mathrm{A}[v, v^{-1}]\) and the A-algebra \(\mathrm{P}(\mathfrak{a}) = \sum_{i \geqslant 0} \mathfrak{a}^{i} v^{i} \subset \mathrm{A}[v]\). These are two graded A-algebras having the same total ring of fractions. Show that for an element \(a \in A\) the following conditions are equivalent:
\end{enumerate}

\begin{enumerate}
\def\labelenumi{(\roman{enumi})}
\item
  The element \(av\) of the \(\mathcal{A}(\mathfrak{a})\)-algebra \(\mathbf{A}[v, v^{-1}]\) is integral over \(\mathcal{A}(\mathfrak{a})\).
\item
  The element av of the P(a)-algebra A{[}v{]} is integral over P(a).
\item
  There exists an integral dependence relation:
\end{enumerate}

\[
a ^ {k} + b _ {1} a ^ {k - 1} + \dots + b _ {k} = 0 \quad \text { with } \quad b _ {j} \in \mathfrak {a} ^ {j}.
\]

One says that \(a\) is integral over the ideal \(\mathfrak{a}\).

\begin{enumerate}
\def\labelenumi{\arabic{enumi})}
\setcounter{enumi}{8}
\tightlist
\item
  \begin{enumerate}
  \def\labelenumii{\alph{enumii})}
  \tightlist
  \item
    Let A and a be as above, and let b be an ideal of A containing a. Show that the following conditions are equivalent:
  \end{enumerate}
\end{enumerate}

\begin{enumerate}
\def\labelenumi{(\roman{enumi})}
\item
  Every element of b is integral over a.
\item
  The \(\mathcal{A}(\mathfrak{a})\)-graded algebra \(\mathcal{A}(\mathfrak{b})\) is integral.
\item
  The graded \(P(\mathfrak{a})\)-algebra \(P(\mathfrak{b})\) is integral.
\item
  There exists an integer \(k \geqslant 0\) such that \(\mathfrak{a}\mathfrak{b}^{k} = \mathfrak{b}^{k+1}\) and therefore \(\mathfrak{a}^{n}\mathfrak{b}^{k} = \mathfrak{b}^{k+n}\) for all positive integers \(n\). We say that \(\mathfrak{b}\) is integral over \(\mathfrak{a}\).
\end{enumerate}

\begin{enumerate}
\def\labelenumi{\alph{enumi})}
\setcounter{enumi}{1}
\item
  Show that there exists a largest ideal among those that are integral over a. It is denoted \(\overline{a}\) and is called the integral closure of a in A. Verify that if \(a \subset b\), then \(\overline{a} \subset \overline{b}\) and that \(\overline{\overline{a}} = \overline{a}\). Show that if the integral closure \(\overline{P(a)}\) is a finitely generated \(P(a)\)-module, then \(\mathfrak{a}\,\overline{\mathfrak{a}^{k}} = \overline{\mathfrak{a}^{k+1}}\) for k sufficiently large.
\item
  Show that if \(a \in A\) satisfies the integral dependence relation
\end{enumerate}

\[
a ^ {k} + b _ {1} a ^ {k - 1} + \dots + b _ {k} = 0 \quad \text { with } \quad k \geqslant 1 \quad \text { and } \quad b _ {i} \in \mathfrak {a} ^ {i},
\]

then \(\mathfrak{a}(\mathfrak{a} + \mathbf{A}a)^{k-1} = (\mathfrak{a} + \mathbf{A}a)^k\). Deduce from this that the integral closure of \(\mathfrak{a}\) in A coincides with the set of elements of \(\mathbf{A}\) which are integral over \(\mathfrak{a}\).

10) Let A be a ring.

\begin{enumerate}
\def\labelenumi{\alph{enumi})}
\item
  Let B be an integral A-algebra. Show that for every element a of A and every ideal \(\mathfrak{a}\) of A, the element a is integral over \(\mathfrak{a}\) if and only if a is integral over the ideal \(\mathfrak{a}B\) generated by \(\mathfrak{a}\) in B.
\item
  Assuming A is Noetherian and reduced, denote \(\mathfrak{p}_{1},\ldots,\mathfrak{p}_{r}\) the minimal prime ideals of A and consider the natural injection:
\end{enumerate}

\[
\mathrm{A} \rightarrow \prod_ {1 \leqslant i \leqslant r} \mathrm{A} / \mathfrak {p} _ {i}.
\]

Show that it makes \(\prod_{1\leqslant i\leqslant r}\mathrm{A} / \mathfrak{p}_i\) into an integral A-algebra.

\begin{enumerate}
\def\labelenumi{\alph{enumi})}
\setcounter{enumi}{2}
\item
  Deduce from this that, given a Noetherian ring A, for an element a of A to be integral over an ideal \(\mathfrak{a}\), it is necessary and sufficient that, for every minimal prime ideal p of A, the image of a in A/p be integral over the ideal \((\mathfrak{a} + p)/p\).
\item
  Show that if \(\overline{a} = \overline{b}\), the ideals \(a\) and \(b\) have the same minimal associated prime ideals.
\end{enumerate}

\begin{enumerate}
\def\labelenumi{\arabic{enumi})}
\setcounter{enumi}{10}
\tightlist
\item
  Let A be a Noetherian ring, and let a, b be two ideals of A such that Supp(a) = Supp(b) = Spec(A). Show that the following conditions are equivalent:
\end{enumerate}

\begin{enumerate}
\def\labelenumi{(\roman{enumi})}
\tightlist
\item
  There exists a finitely generated A-module M such that \(\operatorname{Supp}(\mathbf{M}) = \operatorname{Spec}(\mathbf{A})\) and such that \(bM \subset aM\).\\
\item
  We have \(a, b \subset \overline{a}\).
\end{enumerate}

\begin{enumerate}
\def\labelenumi{\arabic{enumi})}
\setcounter{enumi}{11}
\tightlist
\item
  Let A be a Noetherian ring, and let a, b, and c be three ideals of A such that c is contained in the radical of A.
\end{enumerate}

\begin{enumerate}
\def\labelenumi{\alph{enumi})}
\item
  Show that if one has the inclusion \(\overline{b + c.a} \subset \overline{a}\), then we have \(\overline{b} \subset \overline{a}\).
\item
  Hence, if A is local, and if one has \(b \subset a\) and \(\overline{b} = \overline{a}\), there exists at least one ideal \(b_{1} \subset b\) such that
\end{enumerate}

\[
\overline {{{\mathbf {b} _ {1}}}} = \overline {{{\mathbf {a}}}},
\]

\(\beta)\) \(\mathbf{b}_{1}\) is minimal for this property, that is, if \(\mathbf{b}_{1}^{\prime}\subset \mathbf{b}_{1}\) and \(\overline{\mathbf{b}_1'} = \overline{\mathbf{b}_1}\), then \(\mathbf{b}_1' = \mathbf{b}_1\).

\begin{enumerate}
\def\labelenumi{\arabic{enumi})}
\setcounter{enumi}{12}
\item
  Let A be a local ring and let a and b be two ideals of A distinct from A. Show that if a ⊂ b, then one has \(\overline{a} = \overline{b}\) if and only if the graded algebra gr \(_b(A)\) is integral over the graded algebra gr \(_a(A)\) . Deduce that if a = (x \(_1\) , \ldots, x \(_d\) ) ⊂ b, where d = dim(A), we have \(\overline{a} = \overline{b}\) if and only if the initial forms \(\xi_{1}\) , \ldots, \(\xi_{d}\) in gr \(_b(A)\) of the x \(_i\) are of degree 1 and form a maximal secant sequence for gr \(_b(A)\) .
\item
  Let A be a local ring, m its maximal ideal, and q an ideal of A, distinct from A, such that A/q has finite length. Suppose that A contains a subfield k such that A = k + m.
\end{enumerate}

\begin{enumerate}
\def\labelenumi{\alph{enumi})}
\tightlist
\item
  Show that there exist elements \(x_{1}, \ldots, x_{d}\) of q, where \(d = \dim(A)\), with \(x_{i} \in q^{\delta_{i}}\) such that, for every sufficiently large integer v, one has \(q^{v} = x_{1} \cdot q^{v - \delta_{1}} + \cdots + x_{d} \cdot q^{v - \delta_{d}}\).
\item
  If one assumes that the field A/m is infinite, one can find an ideal a generated by d elements of q, where \(d = \dim(A)\), such that \(\overline{a} = \overline{q}\) (One may use § 7, no. 5.)
\end{enumerate}

\begin{enumerate}
\def\labelenumi{\arabic{enumi})}
\setcounter{enumi}{14}
\tightlist
\item
  Let A be a local, Noetherian, complete ring with maximal ideal m. Set \(d = \dim(A)\). Assume that A contains a subfield k such that \(A = k + m\).
\end{enumerate}

Show that one can find \(x_{1},\ldots,x_{d}\) in m such that:

the morphism \(\varphi:k[[X_{1},...,X_{d}]]\to A\) defined by \(\varphi(X_{i})=x_{i}\) makes A an integral \(k[[X_{1},...,X_{d}]]\)-algebra;

\(\beta)\) denoting \(\mathfrak{m}_0 = (X_1, \ldots, X_d)\) the maximal ideal of \(k[[X_1, \ldots, X_d]]\), one has \(\overline{\mathfrak{m}_0A} = \mathfrak{m}\) and therefore \(e_{\mathfrak{m}_0\mathbf{A}}(\mathbf{A}) = e_{\mathfrak{m}}(\mathbf{A})\).

Show that if \(x_{1},\ldots,x_{d}\) satisfy condition \(\alpha\), condition \(\beta\) is satisfied if and only if the initial forms \(\xi_{1},\ldots,\xi_{d}\) of \(x_{i}\) in \(\mathrm{gr}_{\mathfrak{m}}(\mathbf{A})\) form a maximal secant sequence.

\begin{enumerate}
\def\labelenumi{\arabic{enumi})}
\setcounter{enumi}{15}
\item
  Let A be a Noetherian local ring and \(x = (x_1, \ldots, x_d)\) an ideal generated by a maximal secant sequence for A. Show that we have the inequality \(e_x(A) \leqslant \text{long}(A/x)\) and that equality holds only if the homomorphism \((A/x)[T_1, \ldots, T_d] \xrightarrow{\varphi} \text{gr}_x(A)\) defined by \(\varphi(T_i) = \xi_i\), where \(\xi_i\) is the class of \(x_i\) modulo \(x^2\), is an isomorphism (which means that \((x_1, \ldots, x_d)\) is completely secant).
\item
  Let A be a Noetherian local ring possessing a completely secant sequence of length \(d = \dim(A)\), let \(\mathfrak{a}\) be the ideal generated by such a sequence, and let \(f \in A\) be an element integral over \(\mathfrak{a}\) but not belonging to it (for example \(A = k[[X_1, ..., X_d]]\), where k is a field, \(\mathfrak{a} = (X_1^k, ..., X_d^k)\), and \(f = X_1^{k-1}X_2\)). Let \(\mathfrak{a}' = \mathfrak{a} + fA\). Show that we have the inequality \(e_{\mathfrak{a}'}(A) > \text{long}(A/\mathfrak{a}')\).
\item
  Let A be a ring and \(\mathfrak{a}\) an ideal of A. For every element \(a \in \mathbf{A}\), denote by \(\nu_{\mathfrak{a}}(a)\) the upper bound of the set of integers \(\nu \in \mathbf{N}\) such that \(a \in \mathfrak{a}^{\nu}\).
\end{enumerate}

\begin{enumerate}
\def\labelenumi{\alph{enumi})}
\tightlist
\item
  Show that if \(a \notin \bigcap_{k=0}^{\infty} \mathfrak{a}^k\) the sequence \(\left(\frac{\nu_{\mathfrak{a}}(a^k)}{k}\right)\) is convergent. We denote \(\overline{\nu}_{\mathfrak{a}}(a)\) the limit. If
\end{enumerate}

\(a \in \bigcap_{k=0}^{\infty} \mathfrak{a}^{k}\) one sets \(\overline{\nu}_{\mathfrak{a}}(a) = +\infty\).

\begin{enumerate}
\def\labelenumi{\alph{enumi})}
\setcounter{enumi}{1}
\tightlist
\item
  Show that, for all \(a \in \mathbf{A}\), one has \(\overline{\nu}_{\mathfrak{a}}(a) = \overline{\nu}_{\mathfrak{a}}(a)\) and deduce from this that the following conditions are equivalent:
\end{enumerate}

\begin{enumerate}
\def\labelenumi{(\roman{enumi})}
\item
  \(\overline{\nu}_{\mathfrak{a}}(a)\geqslant 1,\)
\item
  \(a \in \overline{\mathfrak{a}}\).
\end{enumerate}

(Study the monotonicity of the sequence considered at a).)

\begin{enumerate}
\def\labelenumi{\arabic{enumi})}
\setcounter{enumi}{18}
\tightlist
\item
  Let A be a Noetherian ring and q an ideal of A such that A/q has finite length and q is contained in the radical of A. Show that we have
\end{enumerate}

\[
e _ {\mathfrak {q}} (\mathrm{A}) = e _ {\overline {{\mathfrak {q}}}} (\mathrm{A}).
\]

\begin{enumerate}
\def\labelenumi{\arabic{enumi})}
\setcounter{enumi}{19}
\tightlist
\item
  Let A be a Noetherian ring, M a finitely generated A-module, and \(q_{1}, \ldots, q_{k}\) ideals of A such that \(M/q_{i}M\) has finite length for every i. We assume \(M \neq 0\) and set \(d = \dim_{\mathbf{A}}(\mathbf{M})\). a) Show that \(M/q_{1}^{n_{1}} \ldots q_{k}^{n_{k}}M\) is of finite length for every k-tuple of positive integers \((n_{1}, \ldots, n_{k})\).
\end{enumerate}

\begin{enumerate}
\def\labelenumi{\alph{enumi})}
\setcounter{enumi}{1}
\tightlist
\item
  Consider the graded ring H of type \(\mathbf{N}^k\) such that
\end{enumerate}

\[
\mathrm{H} _ {n _ {1}, \dots , n _ {k}} = \mathfrak {q} _ {1} ^ {n _ {1}} \mathfrak {q} _ {2} ^ {n _ {2}} \dots \mathfrak {q} _ {k} ^ {n _ {k}} / \mathfrak {q} _ {1} ^ {n _ {1} + 1} \mathfrak {q} _ {2} ^ {n _ {2}} \dots \mathfrak {q} _ {k} ^ {n _ {k}},
\]

and the graded H-module N such that

\[
 \mathrm{N} _ {n _ {1}, \dots , n _ {k}} = \mathrm{q} _ {1} ^ {n _ {1}} \mathrm{q} _ {2} ^ {n _ {2}} \dots \mathrm{q} _ {k} ^ {n _ {k}} \mathrm{M} / \mathrm{q} _ {1} ^ {n _ {1} + 1} \mathrm{q} _ {2} ^ {n _ {2}} \dots \mathrm{q} _ {k} ^ {n _ {k}} \mathrm{M};
\]

show that N is a finitely generated graded H-module generated by \(N_{0,\ldots,0}\) and a finite number of elements whose degrees are vectors of the canonical basis of \(\mathbf{Z}^{k}\).

From this, deduce that there exists a polynomial \(\mathbf{Q}(\mathrm{T}) \in \mathbf{Z}[\mathrm{T}_1,..,\mathrm{T}_k]\) and integers \(s_1 \geqslant 1, ..., s_k \geqslant 1\) such that we have

\[
\sum_ {n \in \mathbb {N} ^ {k}} \operatorname{long} _ {\mathrm{A}} (\mathrm{H} _ {n}). \mathrm{T} ^ {n} = \frac {\mathrm{Q} (\mathrm{T})}{(1 - \mathrm{T} _ {1}) ^ {s _ {1}} \dots (1 - \mathrm{T} _ {k}) ^ {s _ {k}}}.
\]

(We proceed by induction on d.)

\begin{enumerate}
\def\labelenumi{\alph{enumi})}
\setcounter{enumi}{2}
\tightlist
\item
  From this, deduce, by extending to \(\mathbf{Z}((\mathrm{T}_{1},\ldots,\mathrm{T}_{k}))\) the results of § 6, no. 1, the existence of integers \(c(\alpha_{1},\ldots,\alpha_{k})\) for the sequences \((\alpha_{1},\ldots,\alpha_{k})\) such that \(|\alpha| = \sum_{i=1}^{k}\alpha_{i} = d\) such that we have
\end{enumerate}

\[
e _ {\mathfrak {q} _ {1} ^ {\nu_ {1}} \dots \mathfrak {q} _ {k} ^ {\nu_ {k}}} (\mathbf {M}) = \sum_ {| \alpha | = d} \frac {d !}{\alpha_ {1} ! \dots \alpha_ {k} !} c (\alpha_ {1},..., \alpha_ {k}) v _ {1} ^ {\alpha_ {1}}... v _ {k} ^ {\alpha_ {k}}
\]

for \((\mathbf{v}_1,\dots,\mathbf{v}_k)\) in \(\mathbf{N}^k\)

\begin{enumerate}
\def\labelenumi{\alph{enumi})}
\setcounter{enumi}{3}
\tightlist
\item
  Suppose now that A is local, with infinite residue field κ. Show, using the existence of superficial elements for the graded H-module N, that one has
\end{enumerate}

\[
c (\alpha_ {1}, \dots , \alpha_ {k}) = \inf _ {\mathbf {q} _ {\alpha}} (e _ {\mathbf {q} _ {\alpha}} (\mathbf {M}))
\]

where \(q_{\alpha}\) ranges over the set of ideals of A generated by \(\alpha_{1}\) elements of \(q_{1}\), \(\alpha_{2}\) elements of \(q_{2}\), \ldots, \(\alpha_{k}\) elements of \(q_{k}\). We sometimes denote \(e(\mathbf{q}_{1}^{[\alpha_{1}]} + \cdots + \mathbf{q}_{k}^{[\alpha_{k}]}; \mathbf{M})\) by the integer \(c(\alpha_{1}, \ldots, \alpha_{k})\), abbreviated as \(e(\mathbf{q}_{1}^{[\alpha_{1}]} + \cdots + \mathbf{q}_{k}^{[\alpha_{k}]})\) when M = A. e) Prove the equality

\[
e _ {\mathfrak {q} _ {1}, \mathfrak {q} _ {2}} (\mathbf {M}) = \sum_ {i = 1} ^ {d} \binom {d} {i} e (\mathfrak {q} _ {1} ^ {[ i ]} + \mathfrak {q} _ {2} ^ {[ d - i ]}; \mathbf {M}).
\]

\begin{enumerate}
\def\labelenumi{\alph{enumi})}
\setcounter{enumi}{5}
\tightlist
\item
  Show that we have:
\end{enumerate}

\[
e \left(\overline {{{\mathfrak {q}}}} _ {1} ^ {[ \alpha_ {1} ]} + \dots + \overline {{{\mathfrak {q}}}} _ {k} ^ {[ \alpha_ {k} ]}; \mathrm{M}\right) = e \left(\mathfrak {q} _ {1} ^ {[ \alpha_ {1} ]} + \dots + \mathfrak {q} _ {k} ^ {[ \alpha_ {k} ]}; \mathrm{M}\right)
\]

where \(\overline{\mathfrak{q}}_{i}\) denotes the integral closure in A of the ideal \(\mathfrak{q}_{i}\) (cf. exercise 19).

\begin{enumerate}
\def\labelenumi{\arabic{enumi})}
\setcounter{enumi}{20}
\tightlist
\item
  Let A and A' be two Noetherian rings, and let q (resp. q') be an ideal of A (resp. A') contained in the radical of A (resp. A') and such that \(\operatorname{long}(A/q)\) (resp. \(\operatorname{long}(A'/q'))\) is finite. Let Q be the ideal \(q \otimes A' + A \otimes q'\) of the ring \(A \otimes_{\mathbf{Z}} A'\). Show that it is contained in the radical of \(A \otimes_{\mathbf{Z}} A'\) and that one has the equality
\end{enumerate}

\[
e _ {\mathfrak {Q}} (\mathrm{A} \otimes_ {\mathbf {Z}} \mathrm{A} ^ {\prime}) = e _ {\mathfrak {q}} (\mathrm{A}). e _ {\mathfrak {q} ^ {\prime}} (\mathrm{A} ^ {\prime}).
\]

22) * Let A be a Noetherian integral domain, which is local and complete (resp. a finitely generated algebra over a field k), \(a \in A\) a non-invertible and nonzero element. We make the following assumptions:

α) \(\operatorname{Supp}(A/aA)\) has a single minimal element p.

\(\beta)\) We have \(aA_{p} = pA_{p}\).

γ) A/p is an integrally closed ring.

The aim of this exercise is to show that, for every maximal ideal m of A containing a, the ring \(A_{m}\) is integrally closed and that one has p = aA ("Hironaka's lemma"). In this exercise we shall use the catenarity property. This property is satisfied by integral, Noetherian, local, complete rings, as will be seen in Chapter X. It is also satisfied by integral algebras of finite type over a field (§ 2, n° 4, th. 3).

\begin{enumerate}
\def\labelenumi{\alph{enumi})}
\item
  Show that \(A_p\) is a discrete valuation ring (§ 3, n° 1, prop. 1 and VI, § 3, n° 6, prop. 9).
\item
  We make the hypotheses \(\alpha\) and \(\beta\) and suppose A is integrally closed. Show that one has \(\mathfrak{p} = a\mathbf{A}\) (VII, § 1, n° 4, prop. 8).
\item
  Let A' be the integral closure of A. Prove that A' is an integral, Noetherian, local and complete ring (resp. that it is an algebra of finite type over k). (Note that A is Japanese (IX, § 4, n° 2, th. 2), hence A' is integral, Noetherian, semi-local and complete and, consequently, local (III, § 2, n° 13, corollary to prop. 19). In the case where A is an algebra of finite type over k, use V, § 3, n° 2, th. 2.)
\item
  Show that one has A' ⊗A A\_p = A\_p (use a) and V, § 1, n° 5, prop. 16). Deduce that there exists a single prime ideal p' of A' lying over p and that one has A\_p = A'\_p', aA'\_p' = p'A'\_p'.
\item
  Let q be a minimal element of the support of A'/aA'. Show that q is a prime ideal of A' lying over p. (Note that q has height 1 (§ 3, n° 1, prop. 1) and that consequently, by virtue of catenarity, one has \(\dim(A'/q)=\dim(A)-1\) . Deduce that one has \(\dim(A/(q \cap A))=\dim(A)-1\) and consequently that \(q \cap A = p\) .)
\item
  Deduce from e) that \(aA'\) is a prime ideal of \(A'\) and that one has \(aA' = p' = pA'\) (one may use b)).
\end{enumerate}

g) Show that one has \(\mathbf{A}_{\mathfrak{m}}^{\prime} = \mathbf{A}_{\mathfrak{m}}\) for every maximal ideal m of A containing \(a\). (From the inclusions \(\mathbf{A} \subset \mathbf{A}' \subset \mathbf{A}_{\mathfrak{p}}\), deduce the inclusions \(\mathbf{A}/\mathfrak{p} \subset \mathbf{A}'/\mathfrak{p}' \subset \mathbf{A}_{\mathfrak{p}}/\mathfrak{p}\mathbf{A}_{\mathfrak{p}}\). Deduce that \(\mathbf{A}/\mathfrak{p}\) and \(\mathbf{A}'/\mathfrak{p}'\) have the same field of fractions and consequently, according to \(\gamma)\), that \(\mathbf{A}/\mathfrak{p} = \mathbf{A}'/\mathfrak{p}' = \mathbf{A}' \otimes_{\mathbf{A}} (\mathbf{A}/\mathfrak{p})\), the last equality resulting from \(f)\). One therefore has \((\mathbf{A}'/\mathbf{A}) \otimes_{\mathbf{A}} (\mathbf{A}/\mathfrak{p}) = 0\). Conclude.) Then complete the proof of Hironaka's lemma. *

\begin{enumerate}
\def\labelenumi{\arabic{enumi})}
\setcounter{enumi}{22}
\tightlist
\item
  \begin{enumerate}
  \def\labelenumii{\alph{enumii})}
  \tightlist
  \item
    Let A be a Noetherian local ring such that \(e(A) = 1\). Show that A possesses a single minimal prime ideal \(p\) such that \(\dim(A) = \dim(A/p)\). Show moreover that one has \(\text{long}(A_p) = 1\) and \(e(A/p) = 1\) (\S 7, no. 1, remark 4).
  \end{enumerate}
\end{enumerate}

\begin{enumerate}
\def\labelenumi{\alph{enumi})}
\setcounter{enumi}{1}
\tightlist
\item
  Recall (p. 82, exerc. 10) that a local ring is equidimensional if, for every minimal prime ideal q of A, one has \(\dim(A/q) = \dim(A)\), and that it has no embedded prime ideals if the A-module A has no embedded associated prime ideals (IV, § 2, no. 3, remark). * (It will be seen in Chapter X that the completions of integral quotients of Macaulay local rings possess these two properties.) * Let A be a local, Noetherian, equidimensional ring with no embedded prime ideals, such that \(e(A) = 1\). Show that A is integral.
\end{enumerate}

24) * Let A be a Noetherian local ring, \(\widehat{A}\) its completion. The aim of this exercise is to prove the equivalence of the properties:

\begin{enumerate}
\def\labelenumi{(\roman{enumi})}
\item
  A is regular;
\item
  \(\hat{\mathbf{A}}\) is equidimensional and has no embedded prime ideals, and \(e(\mathbf{A}) = 1\) ;
\item
  \(\hat{\mathbf{A}}\) is integral and \(e(\mathbf{A}) = 1\) .
\end{enumerate}

Prove the implications (i) \(\Rightarrow\) (ii) \(\Rightarrow\) (iii). Note that (iii) implies the equality \(e(\hat{\mathbf{A}}) = 1\) and that, to prove (iii) \(\Rightarrow\) (i), it suffices to prove that \(\hat{\mathbf{A}}\) is regular. One may therefore, to prove (iii) \(\Rightarrow\) (i), suppose A complete, which we shall do from now on. We shall first treat the case where the residue field \(\kappa_{\mathrm{A}}\) has infinitely many elements. One proceeds by induction on the dimension of A. Examine the case \(\dim(A) = 0\), and henceforth suppose \(\dim(A) > 0\). There then exists a superficial element \(x \in \mathrm{A}\) (\S 7, no. 5, remark 4). By virtue of the catenarity property (exerc. 22) and of \S 3, no. 1, prop. 1, for every minimal prime ideal p among those containing xA, one has \(\dim(A/p) = \dim(A/xA)\). Since one has \(e(\mathrm{A}/x\mathrm{A}) = 1\) (\S 7, no. 5, th. 1), there exists a single minimal prime ideal p among those containing xA, and one has \(xA_{p} = pA_{p}\), \(e(\mathrm{A}/p) = 1\) (exerc. 23). By the induction hypothesis, A/p is regular, hence integrally closed (\S 5, no. 2, cor. 1 to th. 1). By Hironaka's lemma (exerc. 22), one deduces p = xA. Hence A/xA is regular and since A is integral, A is regular (\S 5, no. 3, prop. 2).

Suppose now that \(\kappa_{A}\) is a finite field. Prove that it suffices to show that the blowing-up A{]}X{[} (IX, Appendix, no. 2) is regular. Note that A and consequently A{]}X{[} are isomorphic to quotients of a regular ring (IX, § 2, no. 5, th. 3), and \(e(A]X[)=1\). By the theory of Macaulay rings (Chapter X), the completion \(\widehat{A]X[}\) is equidimensional and has no embedded prime ideals. Deduce that \(\widehat{A]X[}\) is integral (exerc. 23) and that \(e(\widehat{A]X[})=1\), hence that \(A]X[\) is regular. Conclude. *

\begin{enumerate}
\def\labelenumi{\arabic{enumi})}
\setcounter{enumi}{24}
\tightlist
\item
  Let \(k\) be a field, A a \(k\) -local algebra, localized at a prime ideal of a \(k\) -algebra of finite type. Show that A is regular if and only if A is integral and \(e(A) = 1\) . (One may be inspired by the method followed in exerc. 24 or also note that since A is an integral quotient of a regular ring, the local ring \(\hat{A}\) is equidimensional and has no embedded prime ideals.)
\end{enumerate}

\chapter{Complete Noetherian local rings}

In this chapter, all rings are assumed commutative; algebras are associative, commutative and unital. We denote by 1 \(_{A}\) the identity element of a ring A.

If A is a ring and p a prime ideal of A, we denote by \(\kappa(p)\) the residue field of the local ring \(A_p\) . If the ring A is local, we denote by \(\mathfrak{m}_A\) its maximal ideal, and \(\kappa_A\) or \(\kappa(\mathfrak{m}_A)\) its residue field.

We say that a ring homomorphism \(\rho: A \to B\) is flat (resp. faithfully flat) if it makes B a flat (resp. faithfully flat) A-module. Recall (I, § 3, n° 5, prop. 9) that if A and B are local, \(\rho\) is faithfully flat if and only if it is flat and local.

\section{WITT VECTORS}

Throughout this section, p denotes a prime number.

\subsection{Witt polynomials}

For every integer \(n \geqslant 0\), we call the element \(\Phi_n\) of \(\mathbf{Z}[X_0, ..., X_n]\) the \(n\)-th Witt polynomial, defined by

\[
\Phi_ {n} \left(\mathrm{X} _ {0}, \dots , \mathrm{X} _ {n}\right) = \sum_ {i = 0} ^ {n} p ^ {i} \mathrm{X} _ {i} ^ {p ^ {n - i}} = \mathrm{X} _ {0} ^ {p ^ {n}} + p \mathrm{X} _ {1} ^ {p ^ {n - 1}} + \dots + p ^ {n} \mathrm{X} _ {n}.\tag{1}
\]

One evidently has \(\Phi_{0}=X_{0}\) and the recurrence relations

\begin{enumerate}
\def\labelenumi{(\arabic{enumi})}
\setcounter{enumi}{1}
\tightlist
\item
\end{enumerate}

\[
\Phi_ {n + 1} \left(\mathrm{X} _ {0}, \dots , \mathrm{X} _ {n + 1}\right) = \Phi_ {n} \left(\mathrm{X} _ {0} ^ {p}, \dots , \mathrm{X} _ {n} ^ {p}\right) + p ^ {n + 1} \mathrm{X} _ {n + 1}\tag{3}
\]

\[
\Phi_ {n + 1} \left(\mathrm{X} _ {0}, \dots , \mathrm{X} _ {n + 1}\right) = \mathrm{X} _ {0} ^ {p ^ {n + 1}} + p \Phi_ {n} \left(\mathrm{X} _ {1}, \dots , \mathrm{X} _ {n + 1}\right).
\]

When one assigns \(X_{i}\) the weight \(p^{i}\), the polynomial \(\Phi_{n}\) is isobaric of weight \(p^{n}\) (A, IV, p.3).

PROPOSITION 1. --- Let A be a filtered ring and \((\mathbf{J}_n)_{n\in \mathbb{Z}}\) its filtration. We assume that \(\mathbf{J}_0 = \mathbf{A}\) and \(p\cdot 1_{\mathbf{A}}\in \mathbf{J}_{1}\). Let m and n be integers such that \(m\geqslant 1\) and \(n\geqslant 0\), and let \(a_0,\dots,a_n,b_0,\dots,b_n\) be elements of A.

\begin{enumerate}
\def\labelenumi{\alph{enumi})}
\tightlist
\item
  If \(a_i \equiv b_i \mod J_m\) for \(0 \leqslant i \leqslant n\), then we have
\end{enumerate}

\[
\Phi_ {i} (a _ {0}, \dots , a _ {i}) \equiv \Phi_ {i} (b _ {0}, \dots , b _ {i}) \bmod \mathrm{J} _ {m + i} \quad \text{for} \quad 0 \leqslant i \leqslant n.
\]

\begin{enumerate}
\def\labelenumi{\alph{enumi})}
\setcounter{enumi}{1}
\tightlist
\item
  Suppose that, for every integer \(k \geqslant 1\) and for all \(x \in \mathbf{A}\), the relation \(p x \in \mathbf{J}_{k+1}\) implies \(x \in \mathbf{J}_k\). If one has \(\Phi_i(a_0, ..., a_i) \equiv \Phi_i(b_0, ..., b_i) \mod \mathbf{J}_{m+i}\) for \(0 \leqslant i \leqslant n\), then we have \(a_i \equiv b_i \mod \mathbf{J}_m\) for \(0 \leqslant i \leqslant n\).
\end{enumerate}

Lemma 1. --- If x and y are two elements of A congruent modulo \(J_{m}\), one has

\[
x ^ {p ^ {n}} \equiv y ^ {p ^ {n}} \bmod J _ {m + n}.
\]

By induction on \(n\) , we reduce to the case where \(n=1\) . Let P denote the polynomial \(\sum_{i=0}^{p-1} X^i Y^{p-1-i}\) of Z{[}X, Y{]}. In view of the hypothesis made on \(x\) and \(y\) , one has P(x, y) \(\equiv\) P(x, x) \(\equiv\) p, \(x^{p-1}\) mod. J \(_m\) . But J \(_m\) + p.A \(\subset\) J \(_1\) , hence P(x, y) \(\in\) J \(_1\) . Finally, \(x^p - y^p = (x - y)\) P(x, y) belongs to J \(_m\) J \(_1\) \(\subset\) J \(_{m+1}\) .

Let us prove a) by induction on n. The case n = 0 is immediate. Suppose \(n \geqslant 1\). Under the hypotheses of a), we have

\begin{enumerate}
\def\labelenumi{(\arabic{enumi})}
\setcounter{enumi}{3}
\tightlist
\item
\end{enumerate}

\(a_{i}^{p}\equiv b_{i}^{p}\bmod \mathrm{J}_{m + 1}\quad \text{for}\quad 0\leqslant i\leqslant n - 1\quad \text{by Lemma } 1.\)

\[
\Phi_ {n - 1} (a _ {0} ^ {p}, \dots , a _ {n - 1} ^ {p}) \equiv \Phi_ {n - 1} (b _ {0} ^ {p}, \dots , b _ {n - 1} ^ {p}) \pmod{J _ {m + n}}\tag{5}
\]

by the induction hypothesis applied to the elements \(a_0^p, \ldots, a_{n-1}^p, b_0^p, \ldots, b_{n-1}^p\) of A, and

\[
\Phi_ {n} (a _ {0},..., a _ {n}) - p ^ {n}. a _ {n} \equiv \Phi_ {n} (b _ {0},..., b _ {n}) - p ^ {n}. b _ {n} \pmod{J _ {m + n}}\tag{6}
\]

by formulas (2) and (5). Since \(a_{n}-b_{n}\) belongs to \(\mathbf{J}_{m}\), the element \(p^{n}a_{n}-p^{n}b_{n}\) belongs to \(\mathbf{J}_{m+n}\), and from (6) we deduce the congruence

\[
\Phi_ {n} (a _ {0}, \dots , a _ {n}) \equiv \Phi_ {n} (b _ {0}, \dots , b _ {n}) \pmod{J _ {m + n}},
\]

whence a).

We prove \(b)\) by induction on \(n\). The case \(n=0\) is immediate. Suppose \(n\geqslant1\). Under the hypotheses of \(b)\), we have \(a_{i}\equiv b_{i}\) mod. \(\mathbf{J}_{m}\) for \(0\leqslant i\leqslant n-1\) by the induction hypothesis, and we deduce as before the congruences (4), (5), and (6). But by hypothesis \(\Phi_{n}(a_{0},...,a_{n})\) and \(\Phi_{n}(b_{0},...,b_{n})\) are congruent mod. \(\mathbf{J}_{m+n}\), and we therefore have \(p^{n}(a_{n}-b_{n})\in\mathbf{J}_{m+n}\). Since the relation \(p x\in\mathbf{J}_{k+1}\) implies \(x\in\mathbf{J}_{k}\) for every \(x\in\mathbf{A}\) and every \(k\geqslant1\), one has \(a_{n}-b_{n}\in\mathbf{J}_{m}\), which completes the proof.

\subsection{\texorpdfstring{The maps \(f, v\) and \(\Phi\)}{The maps f, v, and Φ}}

Let A be a ring. Endow \(\mathbf{A}^{\mathbf{N}}\) with the structure of the product ring. Denote by \(f_{\mathbf{A}}\), or simply \(f\), the endomorphism \((a_{n})_{n\in\mathbf{N}}\mapsto(a_{n+1})_{n\in\mathbf{N}}\) of \(\mathbf{A}^{\mathbf{N}}\). Denote by \(v_{\mathbf{A}}\), or simply \(v\), the endomorphism of the additive group underlying \(\mathbf{A}^{\mathbf{N}}\) which sends \((a_{n})_{n\in\mathbb{N}}\) to \((0,p\cdot a_{0},p\cdot a_{1},...)\).

For every integer \(m \geqslant 0\), denote by \(\Phi_{m}\) the map from \(\mathbf{A}^{\mathbf{N}}\) into A which sends \(\boldsymbol{a} = (a_{n})_{n \in \mathbf{N}}\) to \(\Phi_{m}(a_{0}, \ldots, a_{m})\). We denote by \(\Phi_{\mathrm{A}}\), or simply \(\Phi\), the map \(\boldsymbol{a} \mapsto (\Phi_{n}(\boldsymbol{a}))_{n \in \mathbf{N}}\) of \(\mathbf{A}^{\mathbf{N}}\) into itself.

Lemma 2.--- Let A be a ring endowed with an endomorphism σ satisfying σ(a) ≡ a\^{}p mod. p.A for every a ∈ A. Let n ≥ 1 be an integer and let a\_0, \ldots, a\_\{n-1\} be elements of A. Set u\_i = Φ\_i(a\_0, \ldots, a\_i) for 0 ≤ i ≤ n - 1. Let u\_n be an element of A. The following conditions are equivalent :

\begin{enumerate}
\def\labelenumi{\alph{enumi})}
\item
  There exists \(a_{n} \in A\) such that \(u_{n} = \Phi_{n}(a_{0}, \ldots, a_{n})\).
\item
  \(\sigma(u_{n-1}) \equiv u_n \pmod{p^nA}\).
\end{enumerate}

For \(0 \leqslant i \leqslant n - 1\), one has \(\sigma(a_i) \equiv a_i^p \pmod{pA}\). By Proposition 1 of no. 1, applied to the case where \(\mathbf{J}_k = p^kA\) for \(k \in \mathbf{N}\) and where \(m = 1\), we have the congruence

\[
\Phi_ {n - 1} \left(\sigma \left(a _ {0}\right), \dots , \sigma \left(a _ {n - 1}\right)\right) \equiv \Phi_ {n - 1} \left(a _ {0} ^ {p}, \dots , a _ {n - 1} ^ {p}\right) \pmod{p ^ {n}A},\tag{7}
\]

that is,

\[
\sigma (u _ {n - 1}) \equiv \Phi_ {n - 1} (a _ {0} ^ {p},..., a _ {n - 1} ^ {p}) \pmod{p ^ {n}A}.\tag{8}
\]

Now, according to formula (2), the relation \(u_n = \Phi_n(a_0, \ldots, a_n)\) is equivalent to

\[
u _ {n} = \Phi_ {n - 1} (a _ {0} ^ {p},..., a _ {n - 1} ^ {p}) + p ^ {n} a _ {n}.\tag{9}
\]

The lemma follows from this.

PROPOSITION 2. --- Let A be a ring.

\begin{enumerate}
\def\labelenumi{\alph{enumi})}
\item
  If \(p.1_{\mathbf{A}}\) is not a zero divisor in \(\mathbf{A}\) , the map \(\Phi_{\mathbf{A}}\) is injective.
\item
  If \(p.1_{\mathbf{A}}\) is invertible in \(\mathbf{A}\) , the map \(\Phi_{\mathbf{A}}\) is bijective.
\item
  If \(\sigma\) is an endomorphism of the ring A, satisfying \(\sigma(a) \equiv a^p \bmod pA\) for all \(a \in \mathbf{A}\), the image \(\mathbf{A}'\) of \(\Phi_{\mathbf{A}}\) is a subring of \(\mathbf{A}^{\mathbf{N}}\), stable under \(f_{\mathbf{A}}\) and \(v_{\mathbf{A}}\). It is the set of elements \((u_n)_{n \in \mathbf{N}}\) of \(\mathbf{A}^{\mathbf{N}}\) such that \(\sigma(u_n) \equiv u_{n+1} \bmod p^{n+1}A\) for all \(n \in \mathbf{N}\).
\end{enumerate}

If \(\boldsymbol{a} = (a_n)_{n \in \mathbb{N}}\) and \(\boldsymbol{u} = (u_n)_{n \in \mathbb{N}}\) are elements of \(A^{\mathbb{N}}\) , the relation \(\Phi_A(\boldsymbol{a}) = \boldsymbol{u}\) is equivalent, according to formula (2), to the equalities

\[
\left\{ \begin{array}{l} u _ {0} = a _ {0}, \\ u _ {n} = \Phi_ {n - 1} \left(a _ {0} ^ {p}, \dots , a _ {n - 1} ^ {p}\right) + p ^ {n} a _ {n} \quad \text {for all} \quad n \geqslant 1. \end{array} \right.\tag{10}
\]

Let \(\boldsymbol{u} = (u_{n})_{n \in \mathbb{N}}\) be an element of \(A^{\mathbb{N}}\). When \(p1_{A}\) is a non-zero-divisor in A (resp. when \(p1_{A}\) is invertible in A), there exists at most one sequence \((a_{n})_{n \in \mathbb{N}}\) in A (resp. exactly one sequence \((a_{n})_{n \in \mathbb{N}}\) in A) satisfying the equalities (10), whence a) and b).

We prove \(c)\). By Lemma 2, the image \(\mathbf{A}'\) of \(\Phi_{\mathbf{A}}\) is the set of \(\boldsymbol{u} = (u_n)_{n \in \mathbb{N}}\) in \(\mathbf{A}^{\mathbf{N}}\) such that \(\sigma(u_n) \equiv u_{n+1} \bmod p^{n+1}A\) for all \(n \in \mathbb{N}\). It follows immediately that \(\mathbf{A}'\) is a subring of \(\mathbf{A}^{\mathbf{N}}\), stable under \(f_{\mathbf{A}}\) and \(v_{\mathbf{A}}\).

Remark. — Let \(\boldsymbol{a} = (a_n)_{n \in \mathbb{N}}\) and \(\boldsymbol{u} = (u_n)_{n \in \mathbb{N}}\) be elements of \(A^{\mathbf{N}}\) such that \(\boldsymbol{u} = \Phi_{\mathbf{A}}(\boldsymbol{a})\), and let \(m\) be an integer \(\geqslant 0\). From (10), we deduce the following assertions:

If the \(u_n\), for \(0 \leqslant n \leqslant m\), belong to a subring B of A and if, for every \(x \in A\), the relation \(px \in B\) implies \(x \in B\), then the \(a_n\), for \(0 \leqslant n \leqslant m\), belong to B.

If A is equipped with an \(\mathbf{N}\)-grading, if \(p1_{A}\) is a non-zero-divisor in A, if \(d \in \mathbf{N}\), and if \(u_{n}\) is homogeneous of degree \(dp^{n}\) for \(0 \leqslant n \leqslant m\), then \(a_{n}\) is homogeneous of degree \(dp^{n}\) for \(0 \leqslant n \leqslant m\).

\subsection{Construction of polynomials}

Let A be the ring Z{[}X, Y{]} of polynomials with integer coefficients in the two families of indeterminates \(\mathbf{X} = (\mathbf{X}_{n})_{n \in \mathbb{N}}\) and \(\mathbf{Y} = (\mathbf{Y}_{n})_{n \in \mathbb{N}}\). Let \(\theta\) be the endomorphism of A defined by \(\theta(\mathbf{X}_{n}) = \mathbf{X}_{n}^{p}\) and \(\theta(\mathbf{Y}_{n}) = \mathbf{Y}_{n}^{p}\) for every \(n \in \mathbf{N}\). Then p is not a zero divisor in A, and the set of a in A such that \(\theta(a) \equiv a^{p} \mod p\). A is a subring of A containing the X\_n and the Y\_n, hence equal to A.

By prop. 2, a) and c) of no. 2, there exist elements \(\mathbf{S} = (\mathrm{S}_{n})_{n\in\mathbb{N}}\) , \(\mathbf{P} = (\mathrm{P}_{n})_{n\in\mathbb{N}}\) , \(\mathbf{I} = (\mathrm{I}_{n})_{n\in\mathbb{N}}\) and \(\mathbf{F} = (\mathrm{F}_{n})_{n\in\mathbb{N}}\) of \(\mathrm{A}^{\mathbf{N}}\) characterized respectively by the equalities

\[
\left\{ \begin{array}{l} \Phi_ {\mathbf {A}} (\mathbf {S}) = \Phi_ {\mathbf {A}} (\mathbf {X}) + \Phi_ {\mathbf {A}} (\mathbf {Y}) \\ \Phi_ {\mathbf {A}} (\mathbf {P}) = \Phi_ {\mathbf {A}} (\mathbf {X}) \Phi_ {\mathbf {A}} (\mathbf {Y}) \\ \Phi_ {\mathbf {A}} (\mathbf {I}) = - \Phi_ {\mathbf {A}} (\mathbf {X}) \\ \Phi_ {\mathbf {A}} (\mathbf {F}) = f _ {\mathbf {A}} \big (\Phi_ {\mathbf {A}} (\mathbf {X}) \big). \end{array} \right.\tag{11}
\]

The elements \(S_{n}\), \(P_{n}\), \(I_{n}\), and \(F_{n}\) of A are therefore characterized by the following formulas (where \(n\) ranges over \(\mathbf{N}\)):

\begin{enumerate}
\def\labelenumi{(\arabic{enumi})}
\setcounter{enumi}{11}
\tightlist
\item
\end{enumerate}

\[
\Phi_ {n} \left(\mathrm{S} _ {0}, \dots , \mathrm{S} _ {n}\right) = \Phi_ {n} \left(\mathrm{X} _ {0}, \dots , \mathrm{X} _ {n}\right) + \Phi_ {n} \left(\mathrm{Y} _ {0}, \dots , \mathrm{Y} _ {n}\right),\tag{13}
\]

\[
\Phi_ {n} \left(\mathrm{P} _ {0}, \dots , \mathrm{P} _ {n}\right) = \Phi_ {n} \left(\mathrm{X} _ {0}, \dots , \mathrm{X} _ {n}\right) \Phi_ {n} \left(\mathrm{Y} _ {0}, \dots , \mathrm{Y} _ {n}\right),\tag{14}
\]

\[
\Phi_ {n} (\mathrm{I} _ {0}, \dots , \mathrm{I} _ {n}) = - \Phi_ {n} (\mathrm{X} _ {0}, \dots , \mathrm{X} _ {n}),\tag{15}
\]

\[
\Phi_ {n} \left(\mathrm{F} _ {0}, \dots , \mathrm{F} _ {n}\right) = \Phi_ {n + 1} \left(\mathrm{X} _ {0}, \dots , \mathrm{X} _ {n + 1}\right).
\]

Let us assign \(X_{n}\) and \(Y_{n}\) the weight \(p^{n}\) for every \(n \in \mathbf{N}\). From the remark of no. 2, we deduce the following assertions:

\begin{enumerate}
\def\labelenumi{\alph{enumi})}
\item
  We have \(\mathbf{S}_n\in \mathbf{Z}[\mathbf{X}_0,\dots ,\mathbf{X}_n,\mathbf{Y}_0,\dots ,\mathbf{Y}_n]\), and \(\mathbf{S}_n\) is isobaric of weight \(p^n\).
\item
  We have \(\mathbf{P}_n\in \mathbf{Z}[\mathrm{X}_0,\dots ,\mathrm{X}_n,\mathrm{Y}_0,\dots ,\mathrm{Y}_n]\), and \(\mathbf{P}_n\) is isobaric of weight \(p^n\) in each of the families \((\mathrm{X}_0,\dots ,\mathrm{X}_n)\) and \((\mathrm{Y}_0,\dots ,\mathrm{Y}_n)\).
\item
  We have \(\mathbf{I}_n\in \mathbf{Z}[\mathrm{X}_0,\dots ,\mathrm{X}_n]\), and \(\mathbf{I}_n\) is isobaric of weight \(p^n\).
\item
  We have \(\mathbf{F}_n\in \mathbf{Z}[\mathbf{X}_0,\dots ,\mathbf{X}_{n + 1}]\), and \(\mathbf{F}_n\) is isobaric of weight \(p^{n + 1}\).
\end{enumerate}

Formula (2) allows in practice the determination of the polynomials \(S_{n}\) , \(P_{n}\) , \(I_{n}\) and \(F_{n}\) step by step.

Examples. -- 1) We have

\[
\mathbf {S _ {0}} = \mathbf {X _ {0}} + \mathbf {Y _ {0}}
\]

\[
\mathrm{S} _ {1} = \mathrm{X} _ {1} + \mathrm{Y} _ {1} - \sum_ {i = 1} ^ {p - 1} \frac {1}{p} \binom {p} {i} \mathrm{X} _ {0} ^ {i} \mathrm{Y} _ {0} ^ {p - i}.
\]

Additionally, \(S_{n} - X_{n} - Y_{n}\) belongs to the ring \(Z[X_{0}, ..., X_{n-1}, Y_{0}, ..., Y_{n-1}]\) . 2) We have

\[
\begin{array}{l} \mathbf {P} _ {0} = \mathbf {X} _ {0} \mathbf {Y} _ {0} \\ \mathbf {P} _ {1} = p \mathbf {X} _ {1} \mathbf {Y} _ {1} + \mathbf {X} _ {0} ^ {p} \mathbf {Y} _ {1} + \mathbf {X} _ {1} \mathbf {Y} _ {0} ^ {p}. \end{array}
\]

\begin{enumerate}
\def\labelenumi{\arabic{enumi})}
\setcounter{enumi}{2}
\tightlist
\item
  When \(p \neq 2\) , one has \(I_n = -X_n\) . For \(p = 2\) , one has
\end{enumerate}

\[
\mathrm{I} _ {0} = - \mathrm{X} _ {0}
\]

\[
\mathrm{I} _ {1} = - \left(\mathrm{X} _ {0} ^ {2} + \mathrm{X} _ {1}\right)
\]

\[
\mathrm{I} _ {2} = - \mathrm{X} _ {0} ^ {4} - \mathrm{X} _ {0} ^ {2} \mathrm{X} _ {1} - \mathrm{X} _ {1} ^ {2} - \mathrm{X} _ {2}.
\]

\begin{enumerate}
\def\labelenumi{\arabic{enumi})}
\setcounter{enumi}{3}
\tightlist
\item
  We have
\end{enumerate}

\[
\mathbf {F} _ {0} = \mathbf {X} _ {0} ^ {p} + p \mathbf {X} _ {1}
\]

\[
\mathrm{F} _ {1} = \mathrm{X} _ {1} ^ {p} + p \mathrm{X} _ {2} - \sum_ {i = 0} ^ {p - 1} \binom {p} {i} p ^ {p - i - 1} \mathrm{X} _ {0} ^ {p i} \mathrm{X} _ {1} ^ {p - i}.
\]

Since \(\Phi_n(F_0, \ldots, F_n) \equiv \Phi_n(X_0^p, \ldots, X_n^p) \pmod{p^{n+1}A}\) for all \(n \in \mathbf{N}\) (formulas (2) and (15)), it follows from Proposition 1(b) that \(F_n \equiv X_n^p \pmod{pA}\) for all \(n \in \mathbf{N}\).

Remark. --- Let J be the set of integers \(j \geqslant 1\) . For every element \(j\) of J, define the polynomial \(\varphi_j\) of \(\mathbf{Z}[(\mathbf{X}_j)_{j \in \mathbf{J}}]\) by the formula

\[
\varphi_ {j} = \sum_ {d} d X _ {d} ^ {j / d},
\]

where the sum is over the elements of J that divide j. For every integer \(n \geqslant 0\) , one has

\[
\varphi_ {p ^ {n}} = \Phi_ {n} (\mathrm{X} _ {p ^ {0}},..., \mathrm{X} _ {p ^ {n}}).
\]

For every ring A and every element \(m\) of \(\mathbf{J}\), we denote by \(\varphi_{m}\) the map from \(A^{\mathbf{J}}\) into A that sends \((a_{j})_{j\in\mathbf{J}}\) to \(\varphi_{m}((a_{j})_{j\in\mathbf{J}})\); we denote by \(\varphi_{A}\), or simply \(\varphi\), the map from \(A^{\mathbf{J}}\) into itself that sends \(a = (a_{j})_{j\in\mathbf{J}}\) to \((\varphi_{m}(a))_{m\in\mathbf{J}}\).

Let \(\mathcal{A} = \mathbf{Z}[(X_j)_{j\in\mathbf{J}},(Y_j)_{j\in\mathbf{J}}]\) be the ring of polynomials with integer coefficients in the two families of indeterminates \(\mathbf{X} = (X_j)_{j\in\mathbf{J}}\) and \(\mathbf{Y} = (Y_j)_{j\in\mathbf{J}}\). One can show (p. 51, exerc. 34) that there exist in \(\mathcal{A}\) elements

\[
\boldsymbol {s} = (s _ {j}) _ {j \in \mathbf {J}}, \boldsymbol {p} = (p _ {j}) _ {j \in \mathbf {J}} \quad \text{and} \quad \boldsymbol {i} = (i _ {j}) _ {j \in \mathbf {J}},
\]

characterized by the following equalities:

\[
\begin{array}{l} \varphi_ {\mathcal {A}} (s) = \varphi_ {\mathcal {A}} (\mathbf {X}) + \varphi_ {\mathcal {A}} (\mathbf {Y}) \\ \varphi_ {\mathcal {A}} (p) = \varphi_ {\mathcal {A}} (\mathbf {X}) \varphi_ {\mathcal {A}} (\mathbf {Y}) \\ \varphi_ {\mathcal {A}} (i) = - \varphi_ {\mathcal {A}} (\mathbf {X}). \end{array}
\]

\subsection{The ring W(A) of Witt vectors}

Let A be a ring. If \(\boldsymbol{a} = (a_n)_{n \in \mathbb{N}}\) and \(\boldsymbol{b} = (b_n)_{n \in \mathbb{N}}\) are elements of \(\mathbf{A}^{\mathbf{N}}\), we will denote by \(S_{\mathbf{A}}(\boldsymbol{a}, \boldsymbol{b})\) (resp. \(P_{\mathbf{A}}(\boldsymbol{a}, \boldsymbol{b})\), respectively \(I_{\mathbf{A}}(\boldsymbol{a})\)), or simply by \(S(\boldsymbol{a}, \boldsymbol{b})\) (resp. \(P(\boldsymbol{a}, \boldsymbol{b})\), respectively \(I(\boldsymbol{a})\)), the sequence \((S_n(a_0, ..., a_n; b_0, \ldots, b_n))_{n \in \mathbb{N}}\) (resp. \((P_n(a_0, ..., a_n; b_0, \ldots, b_n))_{n \in \mathbb{N}}\), respectively \((I_n(a_0, ..., a_n))_{n \in \mathbb{N}}\)). By substituting \(a_n\) for \(X_n\) and \(b_n\) for \(Y_n\), for all \(n \in \mathbb{N}\), in formulas (12), (13), and (14), we obtain the equalities

\begin{enumerate}
\def\labelenumi{(\arabic{enumi})}
\setcounter{enumi}{15}
\tightlist
\item
\end{enumerate}

\[
\Phi_ {\mathrm{A}} (\mathrm{S} _ {\mathrm{A}} (a, b)) = \Phi_ {\mathrm{A}} (a) + \Phi_ {\mathrm{A}} (b)\tag{17}
\]

\[
\Phi_ {\mathrm{A}} (\mathrm{P} _ {\mathrm{A}} (a, b)) = \Phi_ {\mathrm{A}} (a) \Phi_ {\mathrm{A}} (b)\tag{18}
\]

\[
\Phi_ {\mathbf {A}} \left(\mathrm{I} _ {\mathbf {A}} (\boldsymbol {a})\right) = - \Phi_ {\mathbf {A}} (\boldsymbol {a}).
\]

We shall denote W(A) the set A \(^{N}\) endowed with the composition laws S \(_{A}\) and P \(_{A}\).

Let \(\rho : \mathbf{B} \to \mathbf{A}\) be a ring homomorphism. We shall denote by \(\rho^{\mathbf{N}}\), or again by \(\mathbf{W}(\rho)\), the map from \(\mathbf{B}^{\mathbf{N}}\) into \(\mathbf{A}^{\mathbf{N}}\) that sends the element \(b = (b_n)_{n \in \mathbb{N}}\) of \(\mathbf{B}^{\mathbf{N}}\) to \((\rho(b_n))_{n \in \mathbb{N}}\). It follows immediately from the definitions that we have

\begin{enumerate}
\def\labelenumi{(\arabic{enumi})}
\setcounter{enumi}{18}
\tightlist
\item
\end{enumerate}

\[
\mathrm{W} (\rho) \circ \mathrm{S} _ {\mathrm{B}} = \mathrm{S} _ {\mathrm{A}} \circ (\mathrm{W} (\rho) \times \mathrm{W} (\rho))\tag{20}
\]

\[
\mathrm{W} (\rho) \circ \mathrm{P} _ {\mathrm{B}} = \mathrm{P} _ {\mathrm{A}} \circ (\mathrm{W} (\rho) \times \mathrm{W} (\rho))\tag{21}
\]

\[
\mathrm{W} (\rho) \circ \mathrm{I} _ {\mathrm{B}} = \mathrm{I} _ {\mathrm{A}} \circ \mathrm{W} (\rho)\tag{22}
\]

\[
\rho^ {\mathbf {N}} \circ \Phi_ {\mathrm{B}} = \Phi_ {\mathrm{A}} \circ \mathrm{W} (\rho).
\]

Lemma 3. --- Let A be a ring. There exists a surjective ring homomorphism \(\rho: \mathbf{B} \to \mathbf{A}\), where \(\mathbf{B}\) is a ring satisfying the following conditions: p is not a zero divisor in \(\mathbf{B}\), and there exists an endomorphism \(\sigma\) of \(\mathbf{B}\) such that \(\sigma(b) \equiv b^p\) mod. \(p\) . \(\mathbf{B}\) for every \(b \in \mathbf{B}\).

Indeed, it suffices to set \(\mathbf{B} = \mathbf{Z}[(X_{a})_{a \in A}]\), to take for \(\sigma\) the endomorphism of \(\mathbf{B}\) defined by \(\sigma(X_{a}) = X_{a}^{p}\) for every \(a \in A\), and to take for \(\rho\) the homomorphism of \(\mathbf{B}\) into A defined by \(\rho(X_{a}) = a\) for every \(a \in A\).

THEOREM 1. --- a) Let A be a (commutative) ring. Equipped with the addition \(S_A\) and the multiplication \(P_A\), \(W(A)\) is a (commutative) ring. The identity element for addition is the sequence \(0_A\), all of whose terms are zero; the identity element for multiplication is the sequence \(1_A\), all of whose terms are zero except the one of index 0, which equals \(1_A\). The additive inverse of an element \(a\) of \(W(A)\) is \(I_A(a)\).

\begin{enumerate}
\def\labelenumi{\alph{enumi})}
\setcounter{enumi}{1}
\item
  Let \(\rho: \mathbf{B} \to \mathbf{A}\) be a ring homomorphism. Then \(\mathbf{W}(\rho): \mathbf{W}(\mathbf{B}) \to \mathbf{W}(\mathbf{A})\) is a ring homomorphism.
\item
  Let A be a ring. The map \(\Phi_{\mathbf{A}}\) is a ring homomorphism from \(W(A)\) into the product ring \(\mathbf{A}^{\mathbf{N}}\). In particular, for every \(n \in \mathbf{N}\), the map \(\Phi_{n}: a \mapsto \Phi_{n}(a_{0}, \ldots, a_{n})\) is a ring homomorphism from W(A) to A.
\end{enumerate}

Given formulas (16), (17), (19), and (20), it suffices to prove assertion \(a\). Let \(\rho: \mathbf{B} \to \mathbf{A}\) be a ring homomorphism satisfying the conditions of Lemma 3. Let \(\mathbf{B}'\) be the subring of \(\mathbf{B}^{\mathbf{N}}\) consisting of the elements \((b_n)_{n \in \mathbf{N}}\) such that \(\sigma(b_n) \equiv b_{n+1} \pmod{p^{n+1}\mathbf{B}}\) for all \(n \in \mathbf{N}\). By Proposition 2 of no. 2, \(\Phi_{\mathbf{B}}\) induces a bijection \(\Phi_{\mathbf{B}}'\) of W(B) onto \(\mathbf{B}'\). In view of formulas (16) to (18) and the relations \(\Phi_n(\mathbf{0}_{\mathbf{B}}) = 0\) and \(\Phi_n(\mathbf{1}_{\mathbf{B}}) = \mathbf{1}_{\mathbf{B}}\) (\(n \in \mathbf{N}\)), one sees by transport of structure that W(B) is a ring, with identity element \(\mathbf{0}_{\mathbf{B}}\) for addition, \(\mathbf{1}_{\mathbf{B}}\) for multiplication, and opposite of b equal to \(\mathbf{I}_{\mathbf{B}}(b)\).

The map W(ρ): W(B) → W(A) is surjective. By formulas (19) and (20), the equivalence relation R on W(B) associated with the map W(ρ) is compatible with the ring structure of W(B). Since W(ρ) induces a bijection Ψ from the quotient ring W(B)/R onto W(A), compatible with the laws of addition and multiplication, assertion a) follows by transport of structure.

DEFINITION 1. --- Let A be a ring. The ring \(W(A)\) is called the ring of Witt vectors with coefficients in A.

For \(\boldsymbol{a}\) in \(W(A)\) and \(n\) in \(\mathbf{N}\), the element \(\Phi_n(\boldsymbol{a}) = \Phi_n(a_0, \ldots, a_n)\) is sometimes called the ghost component of index \(n\) of \(\boldsymbol{a}\).

Remark. --- We take up again the notations of the remark of no. 3. Let A be a ring. If \(a\) and \(b\) are elements of \(\mathbf{A}^{\mathbf{J}}\) and \(r = (r_{j})_{j\in\mathbf{J}}\) is an element of \(\mathcal{A}^{\mathbf{J}}\), denote by \(r_{\mathrm{A}}(a,b)\) the element \((r_{j}(a,b))_{j\in\mathbf{J}}\) of \(\mathbf{A}^{\mathbf{J}}\). Denote by \(U(A)\) the set \(\mathbf{A}^{\mathbf{J}}\) endowed with the composition laws \(s_{\mathrm{A}}\) and \(p_{\mathrm{A}}\). One can show (p. 52, exerc. 35) that, endowed with the addition \(s_{\mathrm{A}}\) and the multiplication \(p_{\mathrm{A}}\), \(U(A)\) is a (commutative) ring; it is called the universal Witt ring of A. The identity element for addition is the element of \(U(A)\) all of whose components are zero; the identity element for multiplication is the element of \(U(A)\) all of whose components are zero except the one of index 1, which equals \(1_{\mathrm{A}}\); the opposite of an element \(a\) of \(U(A)\) is \(i_{\mathrm{A}}(a)\). The map \(\varphi_{\mathrm{A}}\) is a ring homomorphism from \(U(A)\) to the product ring \(\mathbf{A}^{\mathbf{J}}\).

Let \(\rho: \mathbf{B} \to \mathbf{A}\) be a ring homomorphism; denote by \(\mathrm{U}(\rho)\) the map from \(\mathbf{B}^{\mathbf{J}}\) into \(\mathbf{A}^{\mathbf{J}}\) that sends the element \((b_j)_{j \in \mathbf{J}}\) of \(\mathbf{B}^{\mathbf{J}}\) to the element \((\rho(b_j))_{j \in \mathbf{J}}\) of \(\mathbf{A}^{\mathbf{J}}\). One can show (loc. cit.) that \(\mathrm{U}(\rho)\) is a ring homomorphism from \(\mathrm{U}(\mathrm{B})\) into \(\mathrm{U}(\mathrm{A})\).
